% =============================================================================
% SynPred 深度技术分析合集
% Compile: xelatex synpred_analysis.tex (推荐 XeLaTeX,中文 + 数学公式)
% 或:     latexmk -xelatex synpred_analysis.tex
% =============================================================================
\documentclass[11pt,a4paper]{article}

% ---------- 中文支持 (XeLaTeX) ----------
\usepackage{xeCJK}
\setCJKmainfont{PingFang SC}            % macOS 默认中文字体;Linux 改 Noto Serif CJK SC
\setCJKsansfont{PingFang SC}
\setCJKmonofont{PingFang SC}
\setmainfont{TeX Gyre Termes}           % 英文衬线
\setmonofont{Menlo}[Scale=0.85]         % 等宽

% ---------- 数学 / 排版 ----------
\usepackage{amsmath,amssymb,amsthm}
\usepackage{mathtools}
\usepackage{bm}
\usepackage[a4paper,margin=2.2cm]{geometry}
\usepackage{microtype}
\usepackage{enumitem}
\setlist[itemize]{leftmargin=1.4em,topsep=2pt,itemsep=2pt}
\setlist[enumerate]{leftmargin=1.6em,topsep=2pt,itemsep=2pt}

% ---------- 颜色 / 链接 ----------
\usepackage{xcolor}
\definecolor{linkcol}{HTML}{1A5FB4}
\definecolor{citecol}{HTML}{2EC27E}
\definecolor{urlcol}{HTML}{C64600}
\definecolor{codebg}{HTML}{F6F8FA}
\definecolor{codeborder}{HTML}{D0D7DE}
\definecolor{kwcol}{HTML}{CF222E}
\definecolor{strcol}{HTML}{0A3069}
\definecolor{cmtcol}{HTML}{6E7781}

\usepackage[colorlinks=true,
            linkcolor=linkcol,
            citecolor=citecol,
            urlcolor=urlcol,
            bookmarksopen=true]{hyperref}

% ---------- 表格 ----------
\usepackage{array}
\usepackage{booktabs}
\usepackage{longtable}
\usepackage{tabularx}
\usepackage{makecell}
\renewcommand{\arraystretch}{1.15}

% ---------- 代码块 ----------
\usepackage{listings}
\lstset{
    basicstyle=\ttfamily\small,
    backgroundcolor=\color{codebg},
    frame=single,
    framerule=0.4pt,
    rulecolor=\color{codeborder},
    keywordstyle=\color{kwcol}\bfseries,
    stringstyle=\color{strcol},
    commentstyle=\color{cmtcol}\itshape,
    showstringspaces=false,
    breaklines=true,
    breakatwhitespace=true,
    tabsize=2,
    aboveskip=6pt,
    belowskip=6pt,
    columns=fullflexible,
    keepspaces=true,
}
\lstdefinelanguage{yaml}{
    keywords={true,false,null,True,False,None},
    keywordstyle=\color{kwcol}\bfseries,
    sensitive=true,
    comment=[l]{\#},
    commentstyle=\color{cmtcol}\itshape,
    morestring=[b]',
    morestring=[b]",
    stringstyle=\color{strcol},
}

% ---------- 章节样式 ----------
\usepackage{titlesec}
\titleformat{\section}{\Large\bfseries\color{linkcol}}{\thesection}{0.6em}{}
\titleformat{\subsection}{\large\bfseries}{\thesubsection}{0.6em}{}
\titleformat{\subsubsection}{\normalsize\bfseries}{\thesubsubsection}{0.6em}{}

% ---------- 提示框 ----------
\usepackage[most]{tcolorbox}
\newtcolorbox{repobox}[1][]{%
    enhanced, breakable,
    colback=blue!3, colframe=linkcol,
    boxrule=0.5pt, arc=2pt,
    left=8pt, right=8pt, top=4pt, bottom=4pt,
    title=#1,
    fonttitle=\bfseries\small\color{white},
    coltitle=white,
    colbacktitle=linkcol,
}
\newtcolorbox{warnbox}[1][]{%
    enhanced, breakable,
    colback=orange!5, colframe=urlcol,
    boxrule=0.5pt, arc=2pt,
    left=8pt, right=8pt, top=4pt, bottom=4pt,
    title=#1,
    fonttitle=\bfseries\small\color{white},
    coltitle=white,
    colbacktitle=urlcol,
}

% ---------- 自定义命令 ----------
\newcommand{\file}[1]{\texttt{\small #1}}
\newcommand{\code}[1]{\texttt{\small #1}}
\newcommand{\param}[1]{\texttt{\small\color{kwcol}#1}}
\newcommand{\stagestep}[1]{\textbf{STEP #1}}
\newcommand{\norm}[1]{\left\lVert#1\right\rVert}

% =============================================================================
\title{\bfseries SynPred 深度技术分析\\
       \large{Stage 04 / Stage35 / Stage 07 三大模块逐字可重复指南}}
\author{}
\date{\today}

\begin{document}

\maketitle

\begin{abstract}
本文档整合了 SynPred(Structure-to-Synthesis Route Recommendation)项目三个核心模块的逐行级技术分析,达到"按本文步骤逐字执行即可复现"的程度:

\begin{itemize}
    \item \textbf{Part I — Stage 04 Stage3 条件预测模型}:从 Mixture Residual Flow 到 LightGBM Quantile Ensemble,覆盖 5 份训练脚本(\textasciitilde 4000 行)+ 1 份 LGBM 推理脚本。
    \item \textbf{Part II — Stage35 联合排序器}:V21 兼容性 ExtraTrees(主力) + V43 Template-Aware Pairwise Ranker + V3 Learned Regressor 三种排序范式,共 6 份脚本(\textasciitilde 2600 行)+ pipeline 编排 + 配置层。
    \item \textbf{Part III — Stage 07 端到端推理流水线}:28 个 step 全景,\file{run\_pipeline.py} 状态机、配置模板、断点续跑、降级行为、可靠性层 17 个子函数(共 \textasciitilde 4000 行)。
\end{itemize}

每一部分均包含:核心思想、模块架构、逐函数解析、关键工程决策、端到端复现指南、验收清单、常见坑、设计哲学与可改进点。

\end{abstract}

\tableofcontents
\newpage

% ============================================================ Part I
% =============================================================================
% Part I — Stage 04 Stage3 条件预测模型
% =============================================================================
\part*{Part I \quad Stage 04 · Stage3 条件预测模型}
\addcontentsline{toc}{part}{Part I — Stage 04 · Stage3 条件预测模型}

\setcounter{section}{0}

\section{背景与契约}

到 Stage 03 (Data) 结束,SynPred 已经把 \texttt{stage3\_v5\_mixed} NPZ 准备好,每个 NPZ 含:

\begin{center}
\small
\begin{tabular}{lll}
\toprule
张量 & 形状 & 含义 \\
\midrule
\code{x} & $(N, F)$ float32 & hybrid 特征(描述子 + GNN 嵌入)\\
\code{y\_set} & $(N, V)$ float32 & 前驱物 multi-hot(main + aux)\\
\code{y\_cond\_continuous} & $(N, 2)$ float32 & (温度, 时间) 标准化后 \\
\code{y\_cond\_continuous\_mask} & $(N, 2)$ float32 & 0/1 mask \\
\code{y\_cond\_discrete} & $(N, 2)$ int64 & (atmosphere\_coarse\_idx, synthesis\_type\_idx) \\
\code{y\_cond\_discrete\_mask} & $(N, 2)$ float32 & 0/1 mask \\
\code{sample\_id} & $(N,)$ object & 字符串 id \\
\bottomrule
\end{tabular}
\end{center}

加上 \file{condition\_schema.json}(含 \code{continuous\_schema=\{col: \{mean, std, median\}\}}、\code{discrete\_schema=\{col: vocab[]\}}),\file{schema.json}(数据集汇总)。

\textbf{Stage 04 Stage3 的契约}:
\begin{enumerate}
    \item \textbf{训练}:从 NPZ 读 (x, y\_set) 与目标,处理缺失,训练后保存 ckpt。
    \item \textbf{推理}:对 (x, y\_set) 给出多套候选条件(每条样本输出 top-k,典型 k=5\textasciitilde 9),进入 Stage35 联合排序。
\end{enumerate}

下文按 \emph{Mixture Flow $\to$ LGBM $\to$ MLP/Linear/MDN} 四节展开。

% --------------------------------------------------------------- §2
\section{Mixture Residual Flow:\file{train\_condition\_mixture\_flow\_mixed.py}(1310 行)}

\subsection{双阶段思路:残差解耦}

模型核心思想:\textbf{先用 baseline 给条件均值,Flow 只学残差}。

\begin{equation}
\hat c = c_{\text{base}}(x, y_{\text{set}}) + \Delta c, \qquad \Delta c \sim p_\theta(\cdot \mid x, y_{\text{set}})
\end{equation}

为什么这样设计?Stage3 的目标范围很广(温度 100--2000\,°C,时间 0.1--1000\,h),如果直接让神经网络预测分布,小数据下容易学崩。先用一个简单 baseline(Ridge/MLP/ExtraTrees 都行)给一个粗略均值,Flow 只学"baseline 没捕到的偏离"——残差范围小、分布更接近 0,Flow 容易学。

类比:Boosting 的 stage-wise residual 思想,但这里残差是\textbf{分布}而非\textbf{点}。

\subsection{Baseline 选择(行 264--542)}

\textbf{两套 baseline 代码并存}:

\paragraph{(a) \code{Stage3BaselineModel}(行 264--309)} 同一份代码里的小 MLP,作为 torch ckpt 保存。
\begin{itemize}
    \item \code{set\_encoder}: $V \to$ \code{set\_proj\_dim}=256
    \item \code{trunk}: $(F + 256) \to$ \code{hidden\_dims}=$[512, 256]$
    \item \code{disc\_heads[i]}: \code{Linear(trunk\_out, $K_i$)}
    \item \code{cont\_head}: \code{Linear(trunk\_out, n\_cont=2)}
\end{itemize}

输出 dict:\code{\{disc\_logits: List[Tensor], cont\_pred: Tensor (B, 2)\}}。

\paragraph{(b) \code{SklearnBaselineAdapter}(行 318--442)} 适配 sklearn 风格 pickle payload,常见 payload 是 \code{\{"cont\_models": [Ridge1, Ridge2], "disc\_models": [LogReg1, LogReg2], "use\_y\_set": True\}},由 \file{train\_baseline\_linear.py} 输出。\code{object.\_\_setattr\_\_} 绕过 \code{nn.Module} 的 submodule 注册——sklearn 模型不是 torch 模块。

\code{load\_pickle\_baseline\_payload}(行 444+) 兼容三种文件:\file{.pkl}/\file{.pt}/\file{.joblib},并自动判断是 dict-payload(走 sklearn adapter)还是 \code{nn.Module}(走 \code{Stage3BaselineModel})。

\begin{repobox}[可重复要点 \#1]
\param{--baseline\_ckpt} 是必填参数。SynPred 推荐先用 \file{train\_baseline\_linear.py} 训一份 Ridge baseline pickle,再以此 baseline 训 Flow。如果 baseline ckpt 是 torch 形式,需要给 \param{--baseline\_hidden\_dims "512,256"} 还原结构。
\end{repobox}

\subsection{上下文编码(\code{ResidualContextEncoder},行 548--579)}

\begin{lstlisting}
y_set ─ MLP_set([V, 256]) ─┐
                            ├── concat / add ── MLP_trunk([F+256→512→256]) ── h (B, 256)
x ──────────────────────────┘
\end{lstlisting}

两种融合模式:
\begin{itemize}
    \item \param{fuse\_mode="concat"}(默认):\code{fused = cat([x, set\_repr])},trunk 输入维 $F + $ \code{set\_proj\_dim}
    \item \param{fuse\_mode="add"}:\code{fused = x\_proj(x) + set\_repr},trunk 输入维 \code{set\_proj\_dim},需要 \code{x\_proj} 把 $F \to 256$
\end{itemize}

\code{MLP}(行 243--258)是带 LayerNorm + SiLU + Dropout 的标准 stack(可选 layernorm)。输出 \code{h} 维度 = trunk 最后一层(默认 256),叫 \code{ctx\_dim},被三个 head 共享。

\subsection{Mixture-of-Gaussians 残差头(行 582--697)}

\subsubsection{模型成员}

\begin{lstlisting}[language=Python]
self.gating    = MLP([ctx_dim, gating_hidden_dim=128, n_components=3])
self.components = ModuleList([
    MLP([ctx_dim, flow_hidden_dim=256, 2*y_cont_dim]) for _ in range(n_components)
])
self.resid_disc_heads = ModuleList([Linear(ctx_dim, K_i) for K_i in disc_class_sizes])
\end{lstlisting}

\param{n\_components} 默认 3(可调 5)。每个 component MLP 输出 \code{2 * y\_cont\_dim = 4} 维(2 维均值 + 2 维 log\_scale)。

\subsubsection{提取分量参数(\code{\_component\_params},行 634--640)}

\begin{lstlisting}[language=Python]
gating_logits = self.gating(context)               # (B, K)
params = stack([comp(context) for comp in components], dim=1)  # (B, K, 2*D)
means = params[..., :D]                             # (B, K, D=2)
log_scales = params[..., D:].clamp(-5.0, 3.0)       # 数值稳定
\end{lstlisting}

\code{log\_scales} clamp 是关键:防止 $\sigma$ 飘到 0(\code{log\_scale=-5} $\Rightarrow \sigma \approx 0.0067$)或爆炸(\code{log\_scale=3} $\Rightarrow \sigma \approx 20.1$)。

\subsubsection{NLL(行 656--669)}

数学上:
\begin{equation}
p(\Delta c \mid h) = \sum_{k=1}^K \pi_k(h)\, \mathcal{N}\!\left(\Delta c;\,\mu_k(h),\,\mathrm{diag}(\sigma_k^2(h))\right)
\end{equation}
\begin{equation}
\mathcal{L}_{\text{cont}} = -\frac{1}{N}\sum_i \log p(\Delta c_i \mid h_i)
\end{equation}

对应代码:
\begin{lstlisting}[language=Python]
full_rows = (mask > 0.5).all(dim=1)             # 关键:只用所有 cont 维都有标签的样本
n_valid = int(full_rows.sum().item())
if n_valid == 0: return residual.sum()*0.0, 0
r = residual[full_rows]
gating_logits, means, log_scales = self._component_params(context[full_rows])
scales = torch.exp(log_scales)

diff = (r[:, None, :] - means) / clamp(scales, 1e-6)            # (B, K, D)
log_prob_comp = -0.5 * (diff**2 + 2*log_scales + log(2*pi)).sum(-1)   # (B, K)
log_mix = log_softmax(gating_logits, -1)                         # (B, K)
log_prob = logsumexp(log_mix + log_prob_comp, -1)                # (B,)
return -log_prob.mean(), n_valid
\end{lstlisting}

\textbf{关键工程决策}:\code{(mask > 0.5).all(dim=1)} ——只对\textbf{所有}连续维都有标签的样本计算 NLL。这避免温度缺失但时间存在(或反过来)时的偏差。代价是丢一部分样本。

\begin{repobox}[可重复要点 \#2]
若你想最大利用每条样本(逐维 mask),改用 \file{train\_condition\_residual\_mdn\_mixed.py:masked\_nll}(\S\ref{sec:mdn-residual}),它把 dim 维 mask 加权进 \code{log\_prob\_dim},允许部分维有标签即参与训练。这两种处理在 SynPred 数据上 MAE 差异 \textasciitilde 3\%,Flow 版本更稳但 MDN 版本利用率更高。
\end{repobox}

\subsubsection{离散头损失(行 753--761)}

\begin{lstlisting}[language=Python]
def multihead_classification_loss(logits_list, targets, class_weight_tensors, head_weights):
    losses = []
    for i, logits in enumerate(logits_list):
        weight = class_weight_tensors[i] if class_weight_tensors else None
        loss_i = F.cross_entropy(logits, targets[:, i], weight=weight)
        losses.append(loss_i * head_weights[i])
    return stack(losses).sum() / max(sum(head_weights), 1e-8)
\end{lstlisting}

数学上:
\begin{equation}
\mathcal{L}_{\text{disc}} = \frac{\sum_j w_j \cdot \mathrm{CE}(\hat y_j, y_j)}{\sum_j w_j}
\end{equation}

\code{class\_weight\_tensors} 来自 \code{build\_class\_weight\_tensor}(行 742--750):
\begin{lstlisting}[language=Python]
counts = bincount(y_disc[:, i], minlength=K_i)
counts = clip(counts, 1.0, None)            # 避免 0 除
w = counts.sum() / counts                   # 频率倒数
w = w / w.mean()                            # 归一化均值=1
\end{lstlisting}

即\textbf{按 $1/\text{freq}$ 加权}。开关 \param{--use\_class\_weights} 控制。SynPred 默认 ON,因为 atmosphere 的类别极不均衡(\code{air\_or\_oxidizing} 占 \textasciitilde 70\%)。

\subsubsection{离散头的 baseline + residual 拼接(行 887)}

\textbf{关键}:Mixture Flow 的离散头不直接给 logits,而是给"残差 logits",最终 logits = baseline\_logits + residual\_logits:
\begin{lstlisting}[language=Python]
final_disc_logits = [b + r for b, r in zip(base_disc_logits, out["resid_disc_logits"])]
loss_disc = multihead_classification_loss(final_disc_logits, y_disc, ...)
\end{lstlisting}

这是把"残差解耦"思想从连续头扩展到离散头——baseline 给 prior logits,Flow 学修正。

\subsection{推理采样(行 678--697)}

\subsubsection{\code{top1\_residual} —— 确定性中位数}

\begin{lstlisting}[language=Python]
gating_logits, means, _ = self._component_params(context)
top_idx = argmax(gating_logits, dim=1)
return means[arange(B), top_idx]               # (B, D)
\end{lstlisting}

取\textbf{门控最高分量的均值}作为 top-1 预测。这是 Stage3 主线 metric \code{top1\_continuous\_mean\_mae\_raw} 的依据。

\subsubsection{\code{sample\_continuous} —— 随机采样 N 条}

\begin{lstlisting}[language=Python]
gating_logits, means, log_scales = self._component_params(context)
probs = softmax(gating_logits, -1)
cat = Categorical(probs=probs)
comp_idx = cat.sample((n_samples,))            # (S, B) 每样本独立采分量

gather_idx = comp_idx[..., None, None].expand(S, B, 1, D)
chosen_means = gather(means_rep, 2, gather_idx).squeeze(2)
chosen_log_scales = gather(log_scales_rep, 2, gather_idx).squeeze(2)

eps = randn_like(chosen_means)
return chosen_means + exp(chosen_log_scales) * eps    # (S, B, D)
\end{lstlisting}

两步采样:
\begin{enumerate}
    \item 按 gating prob 抽 component。
    \item 在该分量上 reparameterize $\mu + \sigma \epsilon$。
\end{enumerate}

\textbf{注意}:每条样本、每次采样都独立采 component,\textbf{不固定为 top-1 分量}——这保留了多模态。

\subsubsection{\code{evaluate\_split} 的 oracle metric(行 898--907)}

\begin{lstlisting}[language=Python]
sample_resid = model.sample_continuous(out["context"], n_gen_samples)   # (S, B, D)
sample_cont = base_cont.unsqueeze(0) + sample_resid                     # (S, B, D)
target = y_cont.unsqueeze(0).expand(S, B, D)
mask = y_mask.unsqueeze(0).expand(S, B, D)
sqerr = ((sample_cont - target)**2) * mask
denom = clamp(mask.sum(2), min=1.0)
per_sample_err = sqerr.sum(2) / denom            # (S, B)
best_idx = argmin(per_sample_err, dim=0)         # (B,)
oracle_cont = sample_cont[best_idx, arange(B)]   # (B, D)
\end{lstlisting}

\code{oracle\_best\_of\_k} = "如果让 oracle 在 K 个采样里挑最好的,误差是多少"——这是上限性能,反映分布质量而非 top-1 accuracy。SynPred 默认 \param{--n\_gen\_samples 8}。

\subsection{标准化与逆变换(行 173--181, 921--924)}

NPZ 里 \code{y\_cond\_continuous} 已经标准化(stage3 v5 mixed 在 03\_data 里做的)。Flow 学的 residual 也是标准化空间的。最终推理要回到原量纲:

\begin{lstlisting}[language=Python]
def inverse_transform_np(values, stats):
    out = values.copy()
    for i, st in enumerate(stats):
        out[:, i] = out[:, i] * st["std"] + st["mean"]
    return out
\end{lstlisting}

\code{cont\_stats} 来自 \code{condition\_schema["continuous\_schema"][col]},在训练时一并保存进 ckpt 的 \code{schema.cont\_stats},推理时从 ckpt 拿。

\subsection{训练范围 clamp(\code{clip\_to\_train\_range\_fn})}

\begin{lstlisting}[language=Python]
def clip_to_train_range_fn(pred, train_min, train_max):
    return clip(pred, train_min[None], train_max[None])
\end{lstlisting}

\param{--clip\_to\_train\_range} 默认开:推理结果 clamp 到训练集观察到的最小/最大原始量纲。这是保守做法——MoG 容易在小样本边缘外推出 5000\,°C 这种荒谬值,clip 让最差情况不至于太离谱。

\code{train\_min/train\_max} 由 \code{collect\_train\_raw\_minmax}(行 799--811)在原始量纲、masked rows 上求 min/max。

\subsection{训练循环(行 1132--1230)}

\begin{lstlisting}[language=Python]
for epoch in 1..epochs:
    for batch in train_loader:
        with torch.no_grad():
            base_out = baseline(x, y_set)
            base_cont = base_out["cont_pred"]
            base_disc_logits = pad_base_disc_logits(...)

        residual_target = y_cont - base_cont
        out = model(x, y_set)
        final_disc_logits = [b + r for b, r in zip(base_disc_logits, resid_disc_logits)]
        loss_disc = multihead_classification_loss(final_disc_logits, y_disc, ...)
        loss_cont, n_valid = model.nll(residual_target, out["context"], y_mask)
        loss = loss_disc + loss_cont
        backward + clip 5.0 + step

    val_metrics = evaluate_split(baseline, model, val_loader, ...)
    score = val_metrics["top1_continuous_mean_mae_raw"]
    early_stopper.step(score)
    if improved: torch.save(...)
\end{lstlisting}

\textbf{默认超参}(行 1014--1042):

\begin{center}
\small
\begin{tabular}{ll}
\toprule
参数 & 默认值 \\
\midrule
\code{hidden\_dims}          & "512,256" \\
\code{set\_proj\_dim}         & 256 \\
\code{flow\_hidden\_dim}      & 256 \\
\code{n\_components}         & 3 \\
\code{gating\_hidden\_dim}    & 128 \\
\code{dropout}              & 0.1 \\
\code{fuse\_mode}            & concat \\
\code{batch\_size}           & 128 \\
\code{epochs}               & 40 \\
\code{patience}             & 8 \\
\code{lr}                   & $10^{-4}$ \\
\code{weight\_decay}         & $10^{-5}$ \\
\code{metric\_name}          & top1\_continuous\_mean\_mae\_raw \\
\code{n\_gen\_samples}        & 8 \\
\code{clip\_to\_train\_range}  & True \\
\code{grad\_clip}            & 5.0 \\
\code{seed}                 & 42 \\
\bottomrule
\end{tabular}
\end{center}

\begin{repobox}[可重复要点 \#3]
\param{metric\_name} 检查会自动识别"loss/mae/rmse/nll"系是否最小化(\code{is\_minimize\_metric},行 728--730),\code{improved} 方向自动反转。这让你切换不同 metric 时不用动 EarlyStopper 代码。
\end{repobox}

\subsection{输出文件}

\begin{lstlisting}
{run_dir}/
├── best_stage3_condition_mixture_flow_mixed.pt  # model_state + config + schema
├── train_log.json                                # 每 epoch 的 train_loss + val_metrics
├── pred_train.csv / pred_val.csv / pred_test.csv  # baseline / top1 / oracle (norm + raw)
└── final_metrics.json                             # train/val/test 三个 split 的全套指标
\end{lstlisting}

\textbf{预测 CSV 列}(\code{save\_predictions},行 814--832):
\begin{lstlisting}
sample_id, true_<disc>, baseline_<disc>, top1_<disc>, oracle_best_of_k_<disc>,
true_<cont>_norm, true_<cont>_raw, mask_<cont>,
baseline_<cont>_norm, baseline_<cont>_raw,
top1_<cont>_norm, top1_<cont>_raw,
oracle_best_of_k_<cont>_norm, oracle_best_of_k_<cont>_raw
\end{lstlisting}

下游可以直接对比 baseline / top1 / oracle 三档,量化 Flow 比 baseline 提升多少。

% --------------------------------------------------------------- §3
\section{LightGBM Quantile Ensemble:\file{train\_lgbm\_quantile\_ensemble.py}(436 行)}

工程上,Mixture Flow 是"美的",但 LGBM 是"稳的"。SynPred 实测:LGBM 在 MAE 上比 Flow 略好(\textasciitilde 5\%),且训练时间从 30min 降到 3min,无需 GPU。配置里 \param{primary\_model=lgbm} 是默认。

\subsection{任务划分}

LGBM 不学一个端到端混合分布,而是把 Stage3 拆成 4 个独立任务:

\begin{center}
\small
\begin{tabular}{lll}
\toprule
任务 & 目标 & 模型 \\
\midrule
温度回归 & 9 个分位数 (0.1, 0.2, \dots, 0.9) & LightGBM \code{objective="quantile"} \\
时间回归 & 9 个分位数 & 同上 \\
气氛分类 & 二分类(oxidizing vs non-oxidizing)& \code{objective="binary"} \\
时间分桶分类 & 三分类(short/medium/long)& \code{objective="multiclass", num\_class=3} \\
\bottomrule
\end{tabular}
\end{center}

每个任务独立训、独立保存。

\subsection{分位数回归}

\subsubsection{Pinball loss}

LightGBM \code{objective="quantile"} 用 pinball loss:
\begin{equation}
\rho_\tau(u) = \max(\tau u,\;(\tau-1)u) =
\begin{cases}
\tau u, & u \ge 0 \\
(\tau-1)u, & u < 0
\end{cases}
\end{equation}

对真实 $y$ 与预测 $\hat y$,损失 $\rho_\tau(y - \hat y)$ 在 $\tau$-分位数处期望最小,即 LightGBM 学到的 $\hat y$ 是条件 $\tau$-分位数 $Q_\tau(y \mid x)$。

9 个分位数 $\Rightarrow$ 9 个独立模型,每个用 \code{params = \{**QUANTILE\_PARAMS, "alpha": q\}}(行 132)。\param{alpha=q} 就是 LightGBM 中的 $\tau$。

\subsubsection{训练参数(行 42--55)}

\begin{lstlisting}[language=Python]
QUANTILE_PARAMS = {
    "objective": "quantile",
    "metric": "quantile",
    "boosting_type": "gbdt",
    "num_leaves": 255,
    "learning_rate": 0.05,
    "feature_fraction": 0.8,
    "bagging_fraction": 0.8,
    "bagging_freq": 5,
    "min_data_in_leaf": 20,
    "seed": 42,
}
\end{lstlisting}

\param{num\_leaves=255} 偏大,让模型有足够 capacity 学复杂分位数曲面。\param{feature\_fraction=0.8} + \param{bagging\_fraction=0.8} + \param{min\_data\_in\_leaf=20} 起防过拟作用。\param{num\_boost\_round=200},early stopping patience=20。

\subsubsection{mask 处理(行 120--126)}

\begin{lstlisting}[language=Python]
mask_train = data["train"]["y_cond_continuous_mask"][:, target_idx] > 0.5
X_train_m, y_train_m = X_train[mask_train], y_train[mask_train]
\end{lstlisting}

\textbf{只用有标签的行训练}——非常直白,不像 Flow 还要做"逐维 mask 加权"。代价是某些样本对温度有标签但对时间无,温度模型用它,时间模型不用,资源利用率次优但实现极简。

\subsubsection{评估(行 149--178)}

\begin{lstlisting}[language=Python]
top1_pred = test_preds[0.5]                        # median 当 top-1
mae = abs(y - top1_pred).mean()
median_ae = median(abs(y - top1_pred))
r2 = r2_score(y, top1_pred)

# Oracle: 9 个 quantile 里挑误差最小的
all_preds = stack([test_preds[q] for q in QUANTILES], 1)
errors = abs(all_preds - y[:, None])
oracle_pred = all_preds[arange(N), errors.argmin(1)]
oracle_mae = abs(y - oracle_pred).mean()

# 阈值内比例
for thresh in [25, 50, 100, 200]:    # 温度
    metrics[f"top1_within_{thresh}"] = (abs_err <= thresh).mean()
\end{lstlisting}

\textbf{关键工程指标}:\code{top1\_within\_50} = "预测温度误差 $\le$ 50\,°C 的比例",这比 MAE 更直接反映"工程上可用"。SynPred 实测:LGBM \code{top1\_within\_100} \textasciitilde 75\%,Flow top1 \textasciitilde 70\%。

\subsection{气氛二分类 + 时间分桶三分类}

\paragraph{气氛二分类(\code{train\_atmosphere\_classifier},行 181--254)}
简化为二分类(\code{oxidizing} vs \code{inert/reducing/vacuum}):
\begin{lstlisting}[language=Python]
y = (df["target_atmosphere_coarse"] == "air_or_oxidizing").astype(int)
valid = df["has_target"] > 0.5
\end{lstlisting}

不用 NPZ,从 \file{task\_views\_dir/atmosphere\_coarse\_\{train,val,test\}.csv} 读。推理时用 \code{threshold=0.6}(不是 0.5)预测正类——这是 macro\_f1 调出来的。

\paragraph{时间分桶三分类(\code{train\_time\_bucket\_classifier},行 257--332)}
\begin{lstlisting}[language=Python]
classes = ["short", "medium", "long"]
y = df["target_time_bucket"].map(class_map)

# 关键:排除可能泄漏的连续时间列
leak_cols = {"target_time_h_log1p", "target_time_h_clean"}
feat_cols = [c for c in df.columns if c not in meta_cols and c not in leak_cols]
\end{lstlisting}

\textbf{反泄漏}:\code{time\_bucket} 是从 \code{target\_time\_h\_clean} 离散化得到的,如果训练特征里含 \code{target\_time\_h\_log1p} 就泄漏了。

\subsection{推理时的 Quantile 优先级(\file{13b\_run\_stage3\_infer\_lgbm\_quantile.py:210})}

这是 LGBM 模块最优雅的推理 trick:

\begin{lstlisting}[language=Python]
# Order: Q0.5, Q0.4, Q0.6, Q0.3, Q0.7, Q0.2, Q0.8, Q0.1, Q0.9
quantile_priority = [0.5, 0.4, 0.6, 0.3, 0.7, 0.2, 0.8, 0.1, 0.9]
k = min(top_k_conditions, 9)

for sample_i in range(N):
    seen_temps = set()
    exported = 0
    for q in quantile_priority:
        temp_val = temp_preds[q][i]
        time_val = time_preds[q][i]

        # clip 物理范围
        temp_val = max(0.0, min(2000.0, temp_val))
        time_val = max(0.0, min(5000.0, time_val))

        # 按 10°C 粒度去重
        temp_rounded = round(temp_val / 10) * 10
        if temp_rounded in seen_temps:
            continue
        seen_temps.add(temp_rounded)

        # score:median 1.0,向两侧线性递减
        score = 1.0 - abs(q - 0.5) * 2.0

        rows.append({...})
        exported += 1
        if exported >= k:
            break
\end{lstlisting}

\textbf{优先级解读}:
\begin{itemize}
    \item 0.5 ← median,中位/最稳
    \item 0.4, 0.6 ← 中等不确定的左右翼
    \item 0.3, 0.7 ← 中等翼
    \item 0.2, 0.8 ← 大翼
    \item 0.1, 0.9 ← 极端翼
\end{itemize}

效果是:\textbf{先取保守(中位)估计,然后向两侧扩张}。这等价于:在不确定性范围内提供"由保守到激进"的一组候选,让下游 Stage35 排序器有多样选择。

\textbf{去重}:\code{temp\_rounded = round(temp/10)*10} 让相邻分位数若预测温差 $< 10$\,°C 就合并。

\textbf{score}:median 1.0,Q0.4/Q0.6 = 0.8,Q0.3/Q0.7 = 0.6, \dots, Q0.1/Q0.9 = 0.2。

\subsection{命令行}

\begin{lstlisting}[language=bash]
# 训练:
python scripts/04_train/stage3/train_lgbm_quantile_ensemble.py \
  --project_root /Users/wyc/SynPred \
  --num_boost_round 200 \
  --early_stopping_rounds 20

# 推理:
python scripts/07_infer/.../13b_run_stage3_infer_lgbm_quantile.py \
  --conditioned_x_csv .../stage3_x_with_parents.csv \
  --schema_json       runs/stage3/lgbm_quantile_ensemble_v2_fulldata/schema.json \
  --temp_models_dir   runs/stage3/lgbm_quantile_ensemble_v2_fulldata/temperature \
  --time_models_dir   runs/stage3/lgbm_quantile_ensemble_v2_fulldata/time \
  --atm_model         runs/stage3/lgbm_atmosphere_classifier_v1/model_atmosphere_binary_final.txt \
  --time_bucket_model runs/stage3/lgbm_time_bucket_classifier_v1/model_time_bucket.txt \
  --output_dir        runs/stage3_infer_lgbm/<view> \
  --top_k_conditions  5
\end{lstlisting}

% --------------------------------------------------------------- §4
\section{MLP / Linear baselines}

\subsection{\file{train\_mlp\_predictor.py}(781 行)}

每个连续/离散头一个独立 sklearn \code{MLPRegressor}/\code{MLPClassifier}(行 395, 437):
\begin{lstlisting}[language=Python]
MLPRegressor(
    hidden_layer_sizes=(256, 128),
    activation="relu",
    solver="adam",
    alpha=1e-4,                    # L2
    learning_rate="adaptive",
    learning_rate_init=1e-3,
    max_iter=200,
    early_stopping=True,
    validation_fraction=0.1,
    n_iter_no_change=10,
)
\end{lstlisting}

特点:
\begin{itemize}
    \item \textbf{per-head 独立训练},不共享 backbone
    \item 用 \code{StandardScaler} 标准化输入
    \item 缺失标签的样本被 mask 掉(行 387--394)
    \item \textbf{稀疏类别守护}(行 428--435):某离散类别 < 2 个样本时关闭 early\_stopping 防止 train/val 分裂时 val 那一类完全空
\end{itemize}

\code{fit\_report}(行 380--455)记录所有 dropped/skipped 头,便于诊断。输出 pickle 可作为 Mixture Flow 的 \param{--baseline\_ckpt}。

\subsection{\file{train\_baseline\_linear.py}(499 行)}

最干净对照——\textbf{Ridge + LogisticRegression 单层}:
\begin{lstlisting}[language=Python]
Ridge(alpha=1.0, random_state=42)                            # 连续头
LogisticRegression(C=1.0, max_iter=1000, solver="lbfgs")     # 离散头
DummyClassifier(strategy="most_frequent")                    # 单类时 fallback
\end{lstlisting}

无标准化(\param{--standardize} 默认 OFF,因为 Stage 03 已经标准化过了)。\textbf{单类目标}(uniq $\le$ 1)用 \code{DummyClassifier} 代替。

\begin{repobox}[可重复要点 \#4 — SynPred 推荐的 baseline 链]
\begin{enumerate}
    \item 先训 \file{train\_baseline\_linear.py} $\to$ 生成 Ridge baseline pickle
    \item 再训 \file{train\_condition\_mixture\_flow\_mixed.py --baseline\_ckpt <Ridge pickle>} $\to$ Flow 学残差
\end{enumerate}
这种 Linear baseline + MoG 残差的组合在 SynPred 数据上比 MLP-baseline + MoG 略稳(LR 不会 overfit baseline 阶段)。
\end{repobox}

\subsection{\file{train\_condition\_residual\_mdn\_mixed.py}:同款架构但纯 MDN}
\label{sec:mdn-residual}

与 Mixture Flow 的差别:

\begin{center}
\small
\begin{tabular}{lll}
\toprule
维度 & Mixture Flow & Residual MDN \\
\midrule
残差解耦 & $\checkmark$(显式 baseline)& $\checkmark$(显式 baseline)\\
Mixture Head & \code{n\_components=3},各自 MLP & \code{MDNHead(n\_mixtures=5)},共享一层 \\
log\_$\sigma$ clamp & $[-5, 3]$ & $[-7, 5]$(更宽)\\
NLL 实现 & "全维都有标签才计算" & "逐维 mask 加权"(\code{masked\_nll})\\
默认 \code{n\_components} & 3 & 5(更多模式)\\
\bottomrule
\end{tabular}
\end{center}

\code{MDNHead} 用三个 Linear:\code{pi (in→K)}、\code{mu (in→K*D)}、\code{log\_sigma (in→K*D)}。\code{masked\_nll}(行 291--320)的精髓:
\begin{lstlisting}[language=Python]
log_prob_dim = -0.5*(z**2 + 2*log_sigma + log(2*pi))    # (B, K, D)
log_prob_dim = log_prob_dim * mask * head_weights        # 逐维 mask
valid_dim_count = (mask * head_weights).sum(-1).clamp_min(1.0)
log_prob_comp = log_prob_dim.sum(-1) / valid_dim_count   # 平均到有效维
log_pi = log_softmax(pi_logits, -1)
log_mix = logsumexp(log_pi + log_prob_comp, -1)
return -log_mix.mean()
\end{lstlisting}

允许"温度有标签但时间没有"的样本也参与训练。

\subsection{模型选择参考}

\begin{center}
\small
\begin{tabular}{lll}
\toprule
数据规模 & 推荐主力 & 理由 \\
\midrule
$<$ 5k 样本 & Linear baseline & 简单稳健,不容易过拟合 \\
5k -- 20k & LGBM Quantile & 工程稳过 \\
20k+ & Mixture Flow & 表达力强,分布更细 \\
任意规模 & LGBM + Flow 双跑 & 候选互补,Stage35 排序自动选 \\
\bottomrule
\end{tabular}
\end{center}

% --------------------------------------------------------------- §5
\section{端到端复现指南}

\subsection{依赖}

\begin{lstlisting}[language=bash]
pip install numpy pandas scikit-learn torch
pip install lightgbm        # LGBM 分支需要
\end{lstlisting}

\subsection{完整管道}

\begin{lstlisting}[language=bash]
PR=/Users/wyc/SynPred

# === STEP 1: 03_data 已经准备好 NPZ + condition_schema ===
ls $PR/data/interim/generative/stage3_condition_dataset/hybrid_mixed_v1/
# 期望: train.npz val.npz test.npz schema.json condition_schema.json export_summary.json

# === STEP 2: 训练 baseline(Ridge,作为 Flow 的输入) ===
python $PR/scripts/04_train/stage3/train_baseline_linear.py \
  --input_dir   $PR/data/interim/generative/stage3_condition_dataset/hybrid_mixed_v1 \
  --output_dir  $PR/runs/stage3/baseline_ridge_v1 \
  --ridge_alpha 1.0 --logreg_C 1.0

# === STEP 3: 训练 Mixture Residual Flow ===
python $PR/scripts/04_train/stage3/mixed/train_condition_mixture_flow_mixed.py \
  --input_dir     $PR/data/interim/generative/stage3_condition_dataset/hybrid_mixed_v1 \
  --baseline_ckpt $PR/runs/stage3/baseline_ridge_v1/best_model.pkl \
  --run_dir       $PR/runs/stage3/mixture_flow_v1 \
  --hidden_dims "512,256" --set_proj_dim 256 \
  --flow_hidden_dim 256 --n_components 3 --gating_hidden_dim 128 \
  --dropout 0.1 --fuse_mode concat \
  --batch_size 128 --epochs 40 --patience 8 \
  --lr 1e-4 --weight_decay 1e-5 \
  --metric_name top1_continuous_mean_mae_raw \
  --use_class_weights --clip_to_train_range \
  --n_gen_samples 8 --grad_clip 5.0 --device cuda --seed 42

# === STEP 4: 训练 LGBM Quantile Ensemble ===
python $PR/scripts/04_train/stage3/train_lgbm_quantile_ensemble.py \
  --project_root $PR --num_boost_round 200 --early_stopping_rounds 20

# === STEP 5: LGBM 推理 ===
python $PR/scripts/07_infer/.../13b_run_stage3_infer_lgbm_quantile.py \
  --conditioned_x_csv $PR/runs/stage3_infer/conditioned_x.csv \
  --schema_json       $PR/runs/stage3/lgbm_quantile_ensemble_v2_fulldata/schema.json \
  --temp_models_dir   $PR/runs/stage3/lgbm_quantile_ensemble_v2_fulldata/temperature \
  --time_models_dir   $PR/runs/stage3/lgbm_quantile_ensemble_v2_fulldata/time \
  --atm_model         $PR/runs/stage3/lgbm_atmosphere_classifier_v1/model_atmosphere_binary_final.txt \
  --time_bucket_model $PR/runs/stage3/lgbm_time_bucket_classifier_v1/model_time_bucket.txt \
  --output_dir        $PR/runs/stage3_infer/lgbm_quantile_v1 \
  --top_k_conditions 5
\end{lstlisting}

\subsection{验收清单}

\begin{center}
\small
\begin{tabularx}{\linewidth}{lXl}
\toprule
验收点 & 命令 & 期望 \\
\midrule
Flow ckpt 可加载 & \code{torch.load(...)['schema']['cont\_stats']} & 有 mean/std/median \\
Flow 训练收敛 & \code{train\_log.json} 末尾 \code{val.top1\_*\_mae\_raw} & $<$ 100\,°C(温度)\\
Oracle vs Top-1 差距 & 相同指标 oracle / top1 比 & 0.6--0.8 \\
LGBM temp top1\_within\_100 & metrics.json & $>$ 0.7 \\
LGBM atm accuracy & atm\_metrics\_final.json & $>$ 0.85 \\
推理 jsonl 格式 & \code{head -1 test\_candidates.jsonl} & 含 stage3\_score 等 \\
\bottomrule
\end{tabularx}
\end{center}

\subsection{常见坑}

\begin{center}
\small
\begin{tabularx}{\linewidth}{lXl}
\toprule
现象 & 原因 & 解 \\
\midrule
\code{KeyError: y\_cond\_continuous} & NPZ 是 v4(无 mixed) & 用 \file{27\_build\_stage3\_condition\_dataset\_v5\_mixed.py} 重新构 \\
Flow loss NaN 或 inf & log\_scale 没 clamp 飞了 & 检查 clamp(-5, 3);若数据极端,降至 (-3, 2) \\
LGBM 找不到 feature & 输入 csv 列与训练时不一致 & 重生成 conditioned\_x\_csv \\
MLP fit 报 single class & 某离散维 train 全是同一类 & 看 fit\_report 的 dropped\_or\_skipped\_heads \\
Flow oracle $\approx$ top1 & n\_gen\_samples 太小 & 调大 \param{--n\_gen\_samples 32} \\
\bottomrule
\end{tabularx}
\end{center}

% --------------------------------------------------------------- §6
\section{设计哲学与可改进点}

\subsection{七个工程模式}

\begin{enumerate}
    \item \textbf{残差解耦}:Flow / MDN 都用"baseline + residual"两阶段。
    \item \textbf{混合类型分头处理}:连续(Mixture Gaussian)+ 离散(每维独立 CE,带类频权重)+ 缺失(mask 全维或逐维)。
    \item \textbf{不重新拟合 standardizer}:Flow / MDN / LGBM 都从 NPZ 里读 \code{condition\_schema.continuous\_schema} 的 mean/std。
    \item \textbf{Quantile Ensemble 优先级 trick}:不靠采样而靠"事先训好的多分位数",推理时按"由稳到激进"序遍历。
    \item \textbf{train-only stats}:连续 mean/std + class weights + train\_min\_raw/max\_raw 都从 train 算。
    \item \textbf{Schema in ckpt}:\code{best\_*.pt} 自带 disc\_class\_sizes、cont\_col\_names、cont\_stats。
    \item \textbf{多模型并存,Stage35 自决}:LGBM + Flow + MLP 都输出统一格式 candidate JSONL。
\end{enumerate}

\subsection{可改进点}

\begin{enumerate}[label=(\alph*)]
    \item Flow 的 \code{n\_components=3} 偏小,5--7 个更合理。
    \item LGBM Quantile 互相独立训练,单调性可能违反。可加 isotonic regression 后处理。
    \item MDN/Flow 的 baseline 两套代码并存,可统一为 sklearn pickle 路线。
    \item 气氛二分类丢失了 inert vs reducing 区分。
    \item Quantile Ensemble 的 score 是粗的线性映射,可考虑用 quantile 间距估不确定度。
    \item Mixture Flow 的离散头若 baseline 全 0 则等价纯监督,需在 ckpt 里记录 flag。
\end{enumerate}

\subsection{与同类工作的比较}

\begin{center}
\small
\begin{tabular}{lcccc}
\toprule
维度 & 本实现 (Mixture Flow) & Diffusion & Normalizing Flow & Quantile Forest \\
\midrule
表达力 & 中(MoG)& 高 & 高 & 中 \\
训练稳定 & 高 & 中 & 中 & 高 \\
采样速度 & 极快 & 慢 & 中 & 快 \\
缺失数据 & 易 & 难 & 难 & 易 \\
可解释性 & 高 & 低 & 中 & 中 \\
工程依赖 & torch & torch + diffusion 库 & torch + flow 库 & LightGBM \\
\bottomrule
\end{tabular}
\end{center}

\newpage

% ============================================================ Part II
% =============================================================================
% Part II — Stage35 联合排序器
% =============================================================================
\part*{Part II \quad Stage35 联合排序器}
\addcontentsline{toc}{part}{Part II — Stage35 联合排序器}

\setcounter{section}{0}

\section{整体架构}

\subsection{一条数据的旅程}

\begin{lstlisting}
                       ┌── 03_data: candidates JSONL ──┐
   stage2 (gflownet)   │                                │
   stage2 (cvae)       │  N_recipe = N_target × N_p × N_c
   stage3 (mixture)    │                                │
   stage3 (lgbm)       │                                ▼
                       └────────► 12_summarize_routes ──► synthesis_routes_readable.csv
                                                          │
                                                          ▼
                                               14_filter_for_display
                                                          │
                                                          ▼  display_csv
                                          ┌─────────────┬──────────────────┬────────────┐
                                          │             │                  │            │
                                  rule_rerank   learned_rerank      v21_rerank   (备用 v3 / v43)
                                          │             │                  │            │
                                          └────────┬────┴──────────────────┘            │
                                                   ▼                                    │
                                  export_final_top_routes ◄─ 优先级 v21 > learned > rule
                                                   │                                    │
                                                   ▼                                    ▼
                                          final_top_routes.csv   v3_learned + v43_template_aware
                                                   │                                    │
                                                   └──────────────► 联合后处理 ◄──────┘
                                                                          ▼
                                                              final_recommended_routes.{md,csv}
\end{lstlisting}

注意:\textbf{Stage35 不重新计算特征},它读的是上游已经标好元素覆盖、警告、QC、\code{stage3\_score} 的"可读路线 CSV",在此基础上加几列 \code{stage35\_*\_score / stage35\_*\_rank} 重排。这种"sort-only + annotate"的设计让排序层完全可替换、可旁路、可叠加。

\subsection{排序链定义(\file{run\_pipeline.py:24-64})}

\begin{lstlisting}[language=Python]
STEP_FUNCS = [
    ...
    # Stage35 sort chain
    ("summarize_routes",          steps_stage35.summarize_routes),
    ("filter_display_routes",     steps_stage35.filter_display_routes),
    ("stage35_rule_rerank",       steps_stage35.stage35_rule_rerank),
    ("stage35_learned_rerank",    steps_stage35.stage35_learned_rerank),
    ("stage35_v21_rerank",        steps_stage35.stage35_v21_rerank),
    ("best_route_per_precursor",  steps_stage35.best_route_per_precursor),
    ("export_final_top_routes",   steps_final.export_final_top_routes),
    ...
]
\end{lstlisting}

每一步都"非破坏":新增 \code{stage35\_<X>\_score / stage35\_<X>\_rank} 列,不动既有列。下游 \code{export\_final\_top\_routes} 依次找最强可用 ranking,作为 \code{final\_route\_rank}。

\subsection{配置层(\file{configs/full\_route\_stage3.yaml})}

\begin{lstlisting}[language=yaml]
stage35:
  top_n: 30
  rule_script:  '{project_root}/.../19_stage35_rule_route_rerank.py'
  learned_script: '{project_root}/.../20_stage35_learned_route_rerank.py'
  learned_model:  '{project_root}/runs/stage35/route_ranker_hybrid_mixed_v1/stage35_route_ranker.joblib'
  v21_script:        '.../route_ranker/07_apply_stage35_route_ranker_v21_hybrid.py'
  v21_model:         '.../runs/stage35/route_ranker_v2_hybrid_mixed_v1/stage35_route_ranker_v2_extratrees.joblib'
  v21_feature_cols:  '.../runs/stage35/route_ranker_v2_hybrid_mixed_v1/stage35_route_ranker_v2_feature_cols.json'
  best_per_precursor_script: '.../19b_select_best_route_per_precursor.py'

stage35_v43:
  enabled: true
  model_path:        '.../runs/stage35/route_ranker_v43_template_aware/stage35_v43_template_pairwise_chemonly_extratrees.joblib'
  feature_cols_json: '.../runs/stage35/route_ranker_v43_template_aware/stage35_v43_template_pairwise_chemonly_feature_cols.json'
  ...
\end{lstlisting}

\begin{repobox}[可重复要点 \#1]
每个排序器\textbf{模型 + 特征列 JSON} 两件套必须同时部署。脱节会导致维度不匹配 sklearn 报错;更糟的是悄悄列错位、给出有毒分数。
\end{repobox}

% --------------------------------------------------------------- §2
\section{V21 兼容性 ExtraTrees Ranker(主力)}

\subsection{训练数据}

V21 训练数据来自 03\_data 的 \file{24\_build\_stage35\_hardneg\_compat\_dataset.py},$N\times 5$ 行:
\begin{itemize}
    \item 1 条 positive:\code{(x, true\_set)},label=1
    \item 2 条 hard\_negative:\code{(x, modelpred\_set)}(同一目标,模型预测错的高分集合),label=0
    \item 2 条 random\_negative:\code{(x, random\_set)}(其他目标的真值集合),label=0
\end{itemize}

\textbf{张量字段}:
\begin{lstlisting}[language=Python]
x_struct  : (5N, F)    结构特征(标准化前)
precursor_y : (5N, V)  候选前驱物 multi-hot
x_joint   : (5N, F+V)  concat([x_struct, precursor_y]),整体在 train 上拟合 mean/std 再标准化
y         : (5N,)      0/1
\end{lstlisting}

\subsection{模型与训练}

\textbf{模型}:\code{sklearn.ensemble.ExtraTreesClassifier}(joblib 序列化)。SynPred 配置:
\begin{lstlisting}[language=Python]
ExtraTreesClassifier(
    n_estimators=600,           # 600 棵树
    max_depth=None,
    min_samples_leaf=1,
    class_weight='balanced',    # 1:4 不均衡
    n_jobs=-1,
    random_state=42,
)
\end{lstlisting}

为什么用 ExtraTrees 而非 GBDT/MLP?三点工程考量:
\begin{enumerate}
    \item Tree-based 对 multi-hot 输入天然友好。
    \item 不需要标准化也能训(但 SynPred 仍标准化,因 \code{x\_struct} 与 \code{precursor\_y} 量纲差异大)。
    \item \code{predict\_proba} 输出 $[0,1]$ 自然成为 ranker 分数。
\end{enumerate}

\subsection{推理}

\begin{lstlisting}[language=Python]
df = pd.read_csv(input_csv)
feature_cols = json.load(open(feature_cols_json))
X = df[feature_cols].astype(float).values
X = (X - mean) / std

clf = joblib.load(model_path)
proba = clf.predict_proba(X)[:, 1]
df["stage35_v2_prob"] = proba
df["stage35_v21_score"] = proba
df = df.sort_values(["sample_id", "stage35_v21_score"], ascending=[True, False])
df["stage35_v21_rank"] = df.groupby("sample_id").cumcount() + 1
\end{lstlisting}

输出 csv 在原列基础上增加 \code{stage35\_v2\_prob}、\code{stage35\_v21\_score}、\code{stage35\_v21\_rank}。

\subsection{V21 在 chain 里的角色}

V21 是 SynPred 实测\textbf{最强单点排序器}。优势:
\begin{itemize}
    \item 直接学"该 (target, set) 配对是否能合成"——pointwise 二分类。
    \item 训练数据带 hard negative,区分易混淆候选。
    \item ExtraTrees 鲁棒。
\end{itemize}

劣势:
\begin{itemize}
    \item 没显式比较两条候选——pointwise,无法学 "a 比 b 好"。
    \item 化学先验弱——只看 multi-hot,看不到"前驱物属于硝酸盐/碳酸盐路线"。
\end{itemize}

V43 就是为了补这两点。

% --------------------------------------------------------------- §3
\section{V43 Template-Aware Pairwise Ranker}

\subsection{思想}

\textit{"把候选路线显式分类到化学模板(nitrate / carbonate / phosphate / oxide / ...),然后学习'同目标内 a 比 b 好' 的 pairwise 偏好。"}

V43 是 SynPred 中最完整的可学习排序器,分四步:
\begin{itemize}
    \item \textbf{01} 给路线打 template 标签(rule-based 特征工程,不改排名)
    \item \textbf{02} 用一个弱 reward 函数 \code{compute\_template\_quality} 构造同目标的 pairwise 偏好对
    \item \textbf{03} 训练对称化 ExtraTreesClassifier(差分输入)
    \item \textbf{04} 全候选池 all-vs-all 推理,聚合 wins/win\_rate/mean\_prob 给最终分
\end{itemize}

\subsection{STEP 01 · Template 特征工程(\file{01\_add\_route\_template\_features.py},344 行)}

\subsubsection{单个前驱物的 type 分类(\code{precursor\_type},行 65--130)}

\begin{lstlisting}[language=Python]
def precursor_type(p: str) -> set[str]:
    s = clean_str(p)
    compact = s.replace(" ", "")
    types = set()
    elems = parse_elements_from_formula(s)

    # Hydrate(水合)
    if "·" in compact or re.search(r"[·.]\d*H2O", compact):
        types.add("hydrate")

    # 阴离子家族
    if "NO3" in compact: types.add("nitrate")
    if "CO3" in compact: types.add("carbonate")
    if "PO4" in compact or "P2O5" in compact: types.add("phosphate")
    if "SO4" in compact: types.add("sulfate")
    if "SO3" in compact: types.add("sulfite")
    if "OH" in compact: types.add("hydroxide")
    if "NH4" in compact: types.add("ammonium")

    # Selenium 家族
    if "SeO2" in compact or "SeO3" in compact:
        types.add("selenite_selenate")
    elif "Se" in elems:
        types.add("selenide_or_elemental_se")

    # Sulfide
    if "S" in elems and not ("SO4" in compact or "SO3" in compact):
        types.add("sulfide_or_elemental_s")

    # Halogen
    if elems & {"F","Cl","Br","I"}:
        types.add("halide_or_elemental_halogen")

    # Oxide(防止与 oxyanion 重叠)
    oxyanion_types = {"nitrate","carbonate","phosphate","sulfate","sulfite","selenite_selenate"}
    if "O" in elems and not (types & oxyanion_types):
        types.add("oxide_or_oxygen_source")

    # Organic
    if "C" in elems and "H" in elems and "carbonate" not in types:
        types.add("organic_like")

    # 单元素
    if re.fullmatch(r"[A-Z][a-z]?", compact):
        types.add("elemental")

    return types
\end{lstlisting}

工程取舍:
\begin{itemize}
    \item \textbf{集合而非单值}:Co(NO$_3$)$_2$·6H$_2$O 同时是 nitrate + hydrate
    \item \textbf{互斥规则}:含 SO$_4$ 时不打 sulfide;含 oxyanion 时不打 oxide
    \item \textbf{保守优先}:H$_2$O$_2$ 不打 hydrate(严格要求 \code{·xH2O} 写法)
\end{itemize}

\subsubsection{路线整体 template(\code{infer\_primary\_template},行 133--175)}

优先级序:\textbf{离散阴离子家族(磷/硫/硫酸/碳/硝)$>$ 卤素族 $>$ 氢氧化物 $>$ 氧化物 $>$ 单元素 $>$ 有机辅助 $>$ 水合辅助}。

\begin{lstlisting}[language=Python]
priority = [
    ("phosphate_route",  "phosphate"),
    ("sulfate_route",    "sulfate"),
    ("sulfite_route",    "sulfite"),
    ("carbonate_route",  "carbonate"),
    ("nitrate_route",    "nitrate"),
    ("selenite_selenate_route", "selenite_selenate"),
    ("selenide_route",   "selenide_or_elemental_se"),
    ("sulfide_route",    "sulfide_or_elemental_s"),
    ("halide_route",     "halide_or_elemental_halogen"),
    ("hydroxide_route",  "hydroxide"),
    ("oxide_route",      "oxide_or_oxygen_source"),
    ("elemental_route",  "elemental"),
    ("organic_assisted_route", "organic_like"),
    ("hydrate_assisted_route", "hydrate"),
]
present = [name for (name, ty) in priority if type_counts[ty] > 0]
primary = present[0]
secondary = ";".join(present[1:])
confidence = type_counts[priority_type] / total_typed
\end{lstlisting}

\subsubsection{输出 28 列特征}

\begin{lstlisting}
route_template_primary, route_template_secondary, route_template_confidence,
route_template_n_types, route_template_type_signature,
route_has_<oxide|nitrate|carbonate|phosphate|sulfate|sulfite|selenide|...|hydrate|ammonium>_template (15个 0/1),
route_template_is_common_solid_state    # 任一离散阴离子或氢氧化物存在
route_template_is_overly_elemental      # 单元素比例 ≥ 50% 且至少 2 个前驱物
route_template_elemental_ratio
route_template_matches_target_anion
\end{lstlisting}

\subsection{STEP 02 · 弱 Reward + Pairwise 数据集}

\subsubsection{弱 reward 函数(\code{compute\_template\_quality},行 55--104)}

完整公式:
\begin{align}
Q &= 4\cdot \text{cov} - 2\cdot \text{miss} - 0.25\cdot \text{extra} - 1.5\cdot \text{foreign} - 0.5\cdot \text{extra\_nontriv} \notag \\
  &\quad - 0.6\cdot \text{warn} - 0.8\cdot \text{warn\_pen} \notag \\
  &\quad + 0.8\cdot \text{tpl\_match} + 0.4\cdot \text{common\_solid} + 0.2\cdot \text{tpl\_conf} - 0.35\cdot \text{overly\_elem} \notag \\
  &\quad + 0.4\cdot \text{v42\_score} + 0.2\cdot \text{v33\_prob} + 0.1\cdot \text{v32\_score}
\end{align}

各项含义:
\begin{itemize}
    \item \textbf{正信号}(权重最大):元素覆盖 $+4$,template 与目标阴离子匹配 $+0.8$,common\_solid\_state $+0.4$,template\_confidence $+0.2$
    \item \textbf{负信号}:元素缺失 $-2$,多余阴离子轻罚 $-0.25$,foreign\_cation(目标里没有的金属)$-1.5$,extra\_nontrivial $-0.5$
    \item \textbf{诊断信号}:警告数 $-0.6$,加权警告 $-0.8$,过度元素化 $-0.35$
    \item \textbf{prior 信号}:v42 $+0.4$,v33 $+0.2$,v32 $+0.1$
\end{itemize}

\begin{repobox}[可重复要点 \#3]
这一 reward \textbf{不是 ground truth}——SynPred 没有"实验合成成功率"的标签,是用一组化学家可解释的启发式规则模拟"哪些路线更可能成功"。这是个\textbf{弱监督}信号。
\end{repobox}

\subsubsection{同目标 pair 构造(\code{build\_pairs\_for\_group},行 120--189)}

\begin{lstlisting}[language=Python]
d_sorted = d.sort_values("route_quality_v43_weak", ascending=False)
high_pool = d_sorted.head(min(10, n))         # top-10 候选
low_pool  = d_sorted.tail(min(15, n))         # bottom-15 候选

for a in high_pool:
    for b in low_pool:
        if a == b: continue
        gap = q[a] - q[b]
        if abs(gap) < min_quality_gap:        # default 0.25
            continue
        better, worse = (a, b) if gap>0 else (b, a)
        rec = {
            "infer_name": ..., "target_group": infer_name, "label": 1,
            "quality_gap": abs(gap),
            ...
        }
        # 关键:diff__ 特征 = better - worse
        for c in numeric_feature_cols(d):
            rec[f"diff__{c}"] = safe_float(better[c]) - safe_float(worse[c])
        rows.append(rec)
        if len(rows) >= max_pairs_per_group: break
\end{lstlisting}

\textbf{设计要点}:
\begin{enumerate}
    \item \param{top10 $\times$ bottom15} 而非 all-vs-all
    \item \param{min\_quality\_gap=0.25} 过滤太接近的 pair
    \item \code{label=1} 永远(对称化在 03 trainer 里做)
    \item \code{diff\_\_<feature>} 差分特征
\end{enumerate}

\subsection{STEP 03 · 对称化训练(\file{03\_train\_v43\_template\_pairwise\_ranker\_chemonly.py})}

\subsubsection{排除已有排序信号(行 26--48)}

\begin{lstlisting}[language=Python]
EXCLUDE_KEYWORDS = [
    "score","prob","rank","wins","losses","win_rate","mean_prob","local_index",
    "stage35_v21","stage35_v3","stage35_v31","stage35_v32","stage35_v33","stage35_v42",
    "route_warning_adjusted_score","route_warning_score",
]
\end{lstlisting}

只用化学/template 差异,不用前一代 ranker 的分数。让 V43 学到\textbf{互补}信号。

\subsubsection{对称化(行 96--104)}

\begin{lstlisting}[language=Python]
X_pos = df[feature_cols].fillna(0.0).astype(float)
y_pos = np.ones(len(X_pos), dtype=int)

X_neg = -X_pos                         # 关键:翻转所有差分
y_neg = np.zeros(len(X_neg), dtype=int)

X = pd.concat([X_pos, X_neg], ignore_index=True)
y = np.concatenate([y_pos, y_neg], axis=0)
groups = pd.concat([df["target_group"], df["target_group"]], ignore_index=True)
\end{lstlisting}

数学上,训练样本变成:
\begin{equation}
\begin{cases}
\phi(\text{better}) - \phi(\text{worse}) & \to 1 \\
\phi(\text{worse}) - \phi(\text{better}) & \to 0
\end{cases}
\end{equation}

学到的 $f(\Delta\phi)$ 自动具有反对称性:$f(-\Delta) \approx 1 - f(\Delta)$。这是经典 RankNet / pairwise SVM 的核心 trick。

\begin{repobox}[可重复要点 \#4]
如果不做对称化,模型可能学到"better 永远在 diff 公式左边"的虚假 pattern。对称化让 X 分布关于原点对称,迫使 f 学真实的化学差异。
\end{repobox}

\subsubsection{GroupShuffleSplit by target(行 106--111)}

\begin{lstlisting}[language=Python]
splitter = GroupShuffleSplit(n_splits=1, test_size=0.25, random_state=42)
train_idx, test_idx = next(splitter.split(X, y, groups=groups))
\end{lstlisting}

\textbf{反泄漏}:同一 target 的所有 pair 必须整体落在 train 或 test。

\subsubsection{模型(行 118--126)}

\begin{lstlisting}[language=Python]
clf = ExtraTreesClassifier(
    n_estimators=600,
    max_depth=12,                    # 比 V21 严格
    min_samples_leaf=2,
    class_weight='balanced',
    random_state=42,
    n_jobs=-1,
)
\end{lstlisting}

\subsection{STEP 04 · All-vs-All 推理(\file{04\_apply\_v43\_template\_pairwise\_ranker\_chemonly.py})}

\subsubsection{全候选 all-vs-all(行 84--128)}

\begin{lstlisting}[language=Python]
for g in unique_groups:
    idx = where(groups == g)[0]
    ng = len(idx)
    for ii in range(ng):
        i = idx[ii]
        batch_rows = []; opponents = []
        xi = X_base.iloc[i]
        for jj in range(ng):
            if ii == jj: continue
            j = idx[jj]
            xj = X_base.iloc[j]
            diff = xi.values - xj.values   # i - j
            batch_rows.append(diff); opponents.append(j)

        X_pair = pd.DataFrame(batch_rows, columns=feature_cols)
        prob = model.predict_proba(X_pair)[:, 1]    # P(i better than j)
        for p, j in zip(prob, opponents):
            prob_sum[i] += p; n_comp[i] += 1
            if p >= 0.5:
                wins[i] += 1; losses[j] += 1
\end{lstlisting}

每个候选 $i$ 与同 target 的 $\text{ng}-1$ 个对手比较,获得 \code{wins[i]}, \code{losses[i]}, \code{mean\_prob[i]}, \code{win\_rate[i]}。

\textbf{复杂度}:$O(N \times \text{ng}^2)$,每个 target ng $\le 30$ → 900 对/target,可接受。

\subsubsection{最终分数(行 134)}

\begin{equation}
\text{score}_i = 0.7 \cdot \text{mean\_prob}_i + 0.3 \cdot \text{win\_rate}_i
\end{equation}

mean\_prob 是连续信号给细排序;win\_rate 抗噪声。

\subsection{V43 安全门(\file{apply\_v43\_safe\_strict\_gate.py})}

把不安全候选(警告高、QC 不过关)的 v43 score 强制归零:

\begin{lstlisting}[language=Python]
unsafe_mask = (
    (df["route_warning_level"] == "major_warning") |
    (df["precursor_qc_level"] == "major_warning") |
    (df["route_recommendation_status"] == "review_required")
)
df.loc[unsafe_mask, "v43_safe_strict_score"] = 0.0
\end{lstlisting}

% --------------------------------------------------------------- §4
\section{V3 Learned Ranker(全局 fallback)}

\subsection{角色}

V21 训练数据需要"硬负例",V43 的 pairwise 数据需要"benchmark 跑过的 final\_md"。这两个都是\textbf{冷启动困难}。V3 Learned Ranker 是 fallback:直接对一份 final 候选 csv,按 rank 序生成弱标签,训一个全局 ExtraTreesRegressor。

\subsection{STEP 1 · 数据集(\file{build\_v3\_learned\_ranker\_dataset.py})}

\subsubsection{弱标签(行 49--79)}

\begin{lstlisting}[language=Python]
df = df.sort_values("v3_joint_rerank_rank")

n = len(df)
n_pos = int(n * 0.25)              # top 25% 为正
n_neg = int(n * 0.35)              # bottom 35% 为负

labels = np.full(n, 0.5)            # 中间 40% 为不确定
labels[:n_pos] = 1.0
labels[max(n-n_neg, n_pos):] = 0.0

# listwise rank target
df["v3_rank_target"] = 1.0 - (rank-1)/(n-1)
\end{lstlisting}

三档标签 (1.0 / 0.5 / 0.0) 给的是分类信号;\code{v3\_rank\_target} 给的是连续 listwise 信号。SynPred 默认用前者。

\subsubsection{特征列}

\begin{lstlisting}[language=Python]
preferred_numeric = [
    "stage35_v21_score", "stage35_v2_prob", "stage3_score",
    "temperature_c", "time_h",
    "element_coverage", "missing_count", "extra_element_penalty",
    "route_confidence_score", "precursor_qc_score", "v3_joint_feature_score",
]
# 加上所有 v3_ 前缀的数值列(除了 leak 的)
\end{lstlisting}

V3 用\textbf{已有 ranker 的输出}作为特征——这意味着 V3 是\textbf{二阶 ranker},学的是"如何组合多个 V21/V2/V3 信号"。

\subsection{STEP 2 · 训练 + 推理(\file{apply\_v3\_learned\_ranker.py})}

\textbf{安全 override}(行 93--100):
\begin{lstlisting}[language=Python]
df["v3_learned_safety_bucket"] = df.route_recommendation_status.map({
    "recommended": 0,
    "recommended_with_validation": 1,
    "review_required": 2,
}).fillna(1)
sort_cols = ["v3_learned_safety_bucket", "v3_learned_ranker_score"]
\end{lstlisting}

排序时\textbf{先按安全桶升序,再按 v3\_score 降序}——学习信号永远不能 override 化学安全检查。

\textbf{三档分数解读}:
\begin{itemize}
    \item $\ge 0.78$ → \code{high\_learned\_score}
    \item $\ge 0.50$ → \code{medium\_learned\_score}
    \item 否则 → \code{low\_learned\_score}
\end{itemize}

% --------------------------------------------------------------- §5
\section{Final Export 与排序链优先级(\file{steps\_final.py})}

\subsection{优先级(行 23--43)}

\begin{lstlisting}[language=Python]
if "stage35_v21_csv" in r.outputs:        # 最高
    src, source = (..._v21_csv, "stage35_v21")
elif "stage35_learned_csv" in r.outputs:
    src, source = (..._learned_csv, "stage35_learned")
elif "stage35_rule_csv" in r.outputs:
    src, source = (..._rule_csv, "stage35_rule")
else:
    src, source = (display_csv, "display_filtered")
\end{lstlisting}

\subsection{优先级解读}

\begin{itemize}
    \item \textbf{V21}:可学习兼容性二分类,信号最强
    \item \textbf{learned}:旧版 learned 二分类,作为兼容性 fallback
    \item \textbf{rule}:启发式规则得分,稳但粗
    \item \textbf{display}:仅按 \code{stage3\_score} 排序的 raw 结果
\end{itemize}

% --------------------------------------------------------------- §6
\section{端到端复现指南}

\subsection{完整训练管道}

\begin{lstlisting}[language=bash]
PR=/Users/wyc/SynPred

# === STEP 1: 训 V21 ExtraTrees ===
python $PR/scripts/.../route_ranker/06_train_stage35_route_ranker_v21_hybrid.py \
  --input_dir   $PR/data/interim/generative/stage35_hardneg/hybrid/relaxed_only/temperature \
  --output_dir  $PR/runs/stage35/route_ranker_v2_hybrid_mixed_v1 \
  --n_estimators 600 --random_state 42

# === STEP 2: 跑一次 pipeline 得到 final_top_routes.csv ===
cd $PR/scripts/07_infer/structure_to_synthesis_route/pipeline
python run_pipeline.py --config configs/full_route_stage3.yaml --task <some_target>

# === STEP 3: V43 Template-Aware Ranker ===
# 3a) 加 template 特征
python $PR/scripts/.../v43_template_aware/01_add_route_template_features.py \
  --input_csv  /path/to/synthesis_routes_stage35_v33_chemonly_reranked.csv \
  --output_csv /path/to/synthesis_routes_stage35_v43_template_features.csv \
  --top_n 30

# 3b) 构 pairwise 数据集
python $PR/scripts/.../v43_template_aware/02_build_v43_template_pairwise_dataset.py \
  --benchmark_name benchmark_30_clean_v2_final_shell_v11_v33 \
  --project_root $PR \
  --output_csv $PR/runs/stage35/route_ranker_v43_template_aware/v43_template_pairwise.csv \
  --max_pairs_per_group 80 --min_quality_gap 0.25

# 3c) 训练
python $PR/scripts/.../v43_template_aware/03_train_v43_template_pairwise_ranker_chemonly.py \
  --input_csv  $PR/runs/.../v43_template_pairwise.csv \
  --output_dir $PR/runs/stage35/route_ranker_v43_template_aware \
  --n_estimators 600 --max_depth 12 --min_samples_leaf 2 --test_size 0.25

# 3d) 应用
python $PR/scripts/.../v43_template_aware/04_apply_v43_template_pairwise_ranker_chemonly.py \
  --input_csv  /path/to/synthesis_routes_stage35_v43_template_features.csv \
  --model_path $PR/runs/.../stage35_v43_template_pairwise_chemonly_extratrees.joblib \
  --feature_cols_json $PR/runs/.../stage35_v43_template_pairwise_chemonly_feature_cols.json \
  --output_csv /path/to/synthesis_routes_stage35_v43_template_chemonly_reranked.csv \
  --top_n 30

# === STEP 4: V3 Learned Ranker ===
python $PR/scripts/.../pipeline/scripts/build_v3_learned_ranker_dataset.py \
  --input_csv  /path/to/final_top_routes.csv \
  --top_positive_frac 0.25 --bottom_negative_frac 0.35

# === STEP 5: 端到端 pipeline ===
python $PR/scripts/.../pipeline/run_pipeline.py \
  --config $PR/scripts/.../pipeline/configs/full_route_stage3.yaml
\end{lstlisting}

\subsection{常见坑}

\begin{center}
\small
\begin{tabularx}{\linewidth}{XXX}
\toprule
现象 & 原因 & 解 \\
\midrule
V21 推理报维度不匹配 & feature\_cols.json 与 model 不配套 & 确认两者来自同一 run \\
V43 训练 test\_acc=1.0 & 没用 GroupShuffleSplit & 检查 \code{groups=df.target\_group} 传入 \\
V43 推理特征对不上 & base\_cols 在 csv 里缺列 & \code{ensure\_numeric} 自动填 0 \\
V3 score 全是 0.5 & label 全是中性 & 调小 \param{top/bottom\_frac} \\
Final source = display\_filtered & V21/learned/rule 都缺脚本 & 检查 \code{r.cfg["stage35"]["v21\_script"]} 路径 \\
\bottomrule
\end{tabularx}
\end{center}

% --------------------------------------------------------------- §7
\section{设计哲学与可改进点}

\subsection{八个工程模式}

\begin{enumerate}
    \item \textbf{Sort + Annotate, never destroy}:每个 ranker 只新增 \code{stage35\_<X>\_score / rank} 列。
    \item \textbf{优先级 fallback}:final export 按 v21 > learned > rule > display 选最强可用。
    \item \textbf{配置驱动}:每个 ranker 的脚本 / 模型 / feature\_cols 三件套写在 yaml。
    \item \textbf{Pointwise + Pairwise + Listwise 共存}:V21 (pointwise) + V43 (pairwise) + V3 (listwise) 三种范式并行,互补。
    \item \textbf{对称化 pairwise}:迫使模型学反对称偏好。
    \item \textbf{GroupShuffleSplit by target}:防止同目标在 train/test 间泄漏。
    \item \textbf{Excluded keywords}:V43 显式排除既有 ranker 的 score/rank/prob 作为特征。
    \item \textbf{Safety override 优先}:化学安全检查 > 学习排序。
\end{enumerate}

\subsection{可改进点}

\begin{enumerate}[label=(\alph*)]
    \item V21 是 pointwise 不能学 "a 和 b 都能合成,但 a 更典型"。
    \item V43 reward 函数手工调,可用 BO 自动调权重。
    \item V43 only-chemistry 排除了 stage35\_*\_score,丢失有价值的预测信号。可训两版集成。
    \item V3 是 regressor 不是 classifier,可考虑 sigmoid + BCE on soft labels。
    \item All-vs-all $O(\text{ng}^2)$,可用 Bradley-Terry 模型参数化压缩。
    \item 没有 ListNet / LambdaRank,listwise loss 对 top-k 性能通常 +2-3 ndcg。
    \item Safety bucket 是硬 override,可让其进 score 而非分桶。
\end{enumerate}

\subsection{与文献的对照}

\begin{center}
\small
\begin{tabular}{lcccccc}
\toprule
维度 & 本实现 & LambdaMART & XGBoost & LightGBM & RankNet \\
\midrule
pointwise & V21 \checkmark & $\times$ & $\times$ & $\times$ & $\times$ \\
pairwise  & V43 \checkmark & \checkmark & \checkmark & \checkmark & \checkmark \\
listwise  & $\times$ & \checkmark NDCG & \checkmark & \checkmark & $\times$ \\
化学先验  & template-aware \checkmark & 无 & 无 & 无 & 无 \\
模型      & ExtraTrees & GBDT & GBDT & GBDT & DNN \\
训练数据  & 弱 reward & 真 click & 真 query & 真 query & 真 query \\
\bottomrule
\end{tabular}
\end{center}

SynPred 在"无真实标签 + 化学领域知识 + 多模型集成"三点上做了独特取舍。把 chemistry 知识(template / element coverage / warnings)显式注入 reward 而非让 ranker 端到端学,牺牲了一些抽象表达力,换来\textbf{可解释、可审计、可手动修正}——这对合成路线推荐场景至关重要。

\newpage

% ============================================================ Part III
% =============================================================================
% Part III — Stage 07 端到端推理流水线
% =============================================================================
\part*{Part III \quad Stage 07 · 端到端推理流水线}
\addcontentsline{toc}{part}{Part III — Stage 07 · 端到端推理流水线}

\setcounter{section}{0}

\section{整体架构}

\subsection{目录布局}

\begin{lstlisting}
{project_root}/
├── data/poscar/<task>/                  # 用户上传的目标 POSCAR
├── data/interim/                        # 离线训练好的所有特征/词表
├── runs/                                # 离线训练好的所有 ckpt
└── outputs/inference/<infer_name>/
    ├── work_dir/                         # 中间产物(可清理)
    │   ├── split/infer.jsonl
    │   ├── infer_structdesc.csv
    │   ├── chgnet/graph_embed/...
    │   ├── graph_embed/...
    │   ├── stage2_*/, stage2_summary/
    │   ├── stage3_*/, stage3_conditioned_x_*.csv
    │   └── (各种 .summary.json)
    └── out_dir/                          # 最终展示用产物
        ├── stage3_condition_predictions_*/
        ├── routes_flow_fallback_retrieval_baseline_element_reranked/
        │   ├── synthesis_routes_*.csv
        │   ├── final_top_routes*.csv
        │   ├── final_recommended_routes.csv  ← 用户拿这个
        │   └── *.md
        └── pipeline_v3_manifest.json
\end{lstlisting}

\subsection{状态机(\code{PipelineRunner})}

\code{PipelineRunner} 把"分散的命令行调用"组织成"带状态的有限状态机":
\begin{itemize}
    \item \code{self.outputs: Dict[str, str]} — 每步产出的关键文件路径
    \item \code{self.degraded\_steps: List[Dict]} — 哪一步降级 + 原因
    \item \code{self.step\_timings: List[Dict]} — 每步耗时(秒)
    \item \code{record\_degradation(step, reason)} — 一旦降级就 \code{[DEGRADED]} 日志 + 写 manifest
    \item \code{require\_file(path) / require\_dir(path)} — 缺文件直接 raise
    \item \code{step\_enabled(name)} — 看 yaml \code{steps:} 块,某步 false 跳过
    \item \code{restore\_existing\_outputs()} — 启动时扫 work\_dir/out\_dir 下既有产物,符合校验就回填进 \code{self.outputs}
\end{itemize}

\begin{repobox}[可重复要点 \#1]
\code{restore\_existing\_outputs} 的 file 校验是关键(\file{runner.py:14-30})—— 一个 0 字节或只有表头的 csv 不会被认为有效,这避免了"上次跑崩了留下半成品 csv,这次以为有"。
\end{repobox}

\subsection{命令行}

\begin{lstlisting}[language=bash]
python run_pipeline.py \
  --config <yaml>            # 必填,流水线配置
  [--start_from <step_name>] # 从指定 step 重启(后续步骤都执行)
  [--only_step <step_name>]  # 只跑这一步(便于调试)
  [--skip_preflight]         # 跳过 STEP 0 的资源自检
  [--infer_name <X>]         # 覆盖 yaml 里的 infer_name
  [--project_root <path>]    # 覆盖 yaml 里的项目根
\end{lstlisting}

\textbf{配置加载}(\file{config.py},76 行):
\begin{lstlisting}[language=Python]
def load_config(path):
    cfg = yaml.safe_load(...)
    for _ in range(5):                  # 迭代 5 次解析模板
        new_cfg = resolve_templates(cfg, cfg)
        if new_cfg == cfg: break
        cfg = new_cfg
    cfg["_config_path"] = str(Path(path).resolve())
    return cfg
\end{lstlisting}

\code{resolve\_templates} 把 \code{\{project\_root\}} / \code{\{infer\_name\}} / \code{\{stage2.gflownet\_run\_dir\}}(支持 \code{.} 路径)替换为实际值。\textbf{5 次迭代}是为了支持模板里嵌模板。

\subsection{28 步全景}

\begin{lstlisting}
STEP 0  preflight                       — 资源/CKPT/词表自检
STEP 1  make_infer_split                — POSCAR → infer.jsonl
STEP 2  build_structdesc                — 结构描述子 CSV
STEP 3  build_chgnet_embedding          — CHGNet 64-dim crystal_fea
STEP 4  finalize_graph_embedding        — 优先 CGCNN,缺失则 CHGNet
STEP 5  build_stage2_features           — Stage2 hybrid CSV
STEP 6  build_stage2_npz                — 转 NPZ + meta(用 train 时 μ,σ)
STEP 7  sample_stage2_gflownet          — standard 或 composition_biased
STEP 8  constrain_stage2_by_composition — 后置硬过滤
STEP 9  summarize_stage2                — unique_sets + count
STEP 10 add_composition_fallback        — 元素覆盖低补"组成回退"前驱物
STEP 11 retrieve_stage2_candidates      — k-NN 检索相似目标的真前驱物集合
STEP 12 predict_stage2_baseline         — ExtraTrees 多标签基线
STEP 13 merge_stage2_sources            — 四源合并
STEP 14 rerank_stage2_by_elements       — 元素覆盖加权重排
STEP 15 fix_stage2_global_rank          — 修补合并后 global rank
STEP 16 build_stage3_features           — Stage3 hybrid CSV
STEP 17 build_stage3_conditioned_table  — (x, parent_set) 行展开
STEP 18  run_stage3_flow                — Mixture Flow 推理(默认关)
STEP 18b run_stage3_lgbm                — LightGBM quantile 推理(默认开,主力)
STEP 19  compare_stage3_models          — 二者一致性报表
STEP 20  summarize_routes               — stage3_flat → 可读路线
STEP 21  filter_display_routes          — 物理上限/下限
STEP 22  stage35_rule_rerank            — 规则打分
STEP 23  stage35_learned_rerank         — ExtraTrees regressor 重排
STEP 24  stage35_v21_rerank             — ExtraTrees pairwise 重排(主力)
STEP 25  best_route_per_precursor       — 同一前驱物只留最佳条件
STEP 26  export_final_top_routes        — 选最强可用 ranking
STEP 27  reliability_layer              — 6+ 子步骤
STEP 28  select_final_recommended_routes — 最终稳定列写出
\end{lstlisting}

% --------------------------------------------------------------- §2
\section{STEP 0--4:结构入口与图嵌入}

\subsection{STEP 0 · \code{preflight}}

不做任何写入,只检查上游 ckpt / 词表 / 模板目录就位:
\begin{lstlisting}[language=Python]
r.require_dir(paths["poscar_dir"])
r.require_file(stage2["gflownet_ckpt"])
r.require_dir(stage2["template_dir"])
r.require_file(template_dir / "feature_cols.json")
r.require_file(template_dir / "feature_mean.npy")     # train 期统计的 μ
r.require_file(template_dir / "feature_std.npy")      # σ
...
\end{lstlisting}

\textbf{为什么 preflight 要单独一步?}因为后续 step 跑到一半发现缺 ckpt 才报错,前面的中间产物白算。preflight 只读元数据,几毫秒,但能阻止 30min 后白跑。

\subsection{STEP 1 · \code{make\_infer\_split}}

\begin{lstlisting}[language=bash]
python pipeline/src/01_make_infer_split_from_poscars.py \
  --poscar_dir <poscar_dir> \
  --output_dir <work_dir>/split
# 产出:<work_dir>/split/infer.jsonl
\end{lstlisting}

\code{infer.jsonl} 每行一条:
\begin{lstlisting}[language=Python]
{"sample_id": "abc", "formula": "Fe2O3", "poscar_path": ".../abc.vasp", ...}
\end{lstlisting}

\subsection{STEP 2 · \code{build\_structdesc}}

\begin{lstlisting}[language=bash]
python pipeline/src/02_build_infer_structdesc_direct.py \
  --infer_jsonl <infer.jsonl> \
  --output_csv <work_dir>/infer_structdesc.csv
\end{lstlisting}

走 Stage 03 同样的描述子提取流程(matminer + magpie + composition + symmetry),无 GNN。输出 \file{infer\_structdesc.csv},每行一条,$\sim 120$ 列描述子。

\subsection{STEP 3 · \code{build\_chgnet\_embedding}}

\begin{lstlisting}[language=bash]
python pipeline/src/03_build_infer_graph_embeddings_chgnet.py \
  --infer_jsonl ... --work_dir <work_dir>/chgnet \
  --project_root ... --precursor_vocab_json ... \
  --train_mode gold_only --max_sites 200
\end{lstlisting}

调用 CHGNet 预训练模型抽 \code{crystal\_fea}(默认 64 维)。对 site 数 $> 200$ 的超晶胞会触发 fallback,通常用对称化后的 primitive cell 替代。

\subsection{STEP 4 · \code{finalize\_graph\_embedding}}

\textbf{双路径}:
\begin{enumerate}
    \item \textbf{Primary}:CGCNN 缓存 + checkpoint + model\_py + model\_class 都齐 $\to$ 跑 04\_finalize 把 CGCNN embedding 与 CHGNet 拼接
    \item \textbf{Fallback}:任一缺失 $\to$ \code{record\_degradation("finalize\_graph\_embedding", "CGCNN unavailable, using CHGNet-only")},直接把 CHGNet csv 抄到 \code{final\_graph\_embed\_csv}
\end{enumerate}

\begin{repobox}[可重复要点 \#3]
这是流水线第一个"显式 degradation"。SynPred 哲学:任何上游模型缺失,流水线\textbf{降级而非崩溃}——下游模型可能仍能工作。但 manifest 必须如实记录降级。
\end{repobox}

% --------------------------------------------------------------- §3
\section{STEP 5--15:Stage2 候选集合的"四源合并"}

\subsection{STEP 5 · \code{build\_stage2\_features}}

\begin{lstlisting}[language=bash]
python pipeline/src/05_build_hybrid_features_infer.py \
  --task stage2 \
  --output_dir <work_dir>/stage2_hybrid_csv \
  --infer_descriptor_csv ... \
  --infer_embedding_csv  ... \
  --embedding_prefix graph_emb \
  --replicate_to_train_val
\end{lstlisting}

\code{--replicate\_to\_train\_val} 是个奇怪但聪明的开关:即使是推理,也写出三份 train/val/test 同样内容,让下游"以为是数据集"。

\subsection{STEP 6 · \code{build\_stage2\_npz}}

\begin{lstlisting}[language=bash]
python pipeline/src/07_build_stage2_gflownet_infer_npz.py \
  --infer_hybrid_csv ... \
  --template_dir <stage2_gflownet_dataset/hybrid/gold_only> \
  --output_dir <work_dir>/stage2_hybrid \
  --split_name test
\end{lstlisting}

\textbf{关键 trick}:\code{--template\_dir} 指向训练时的 NPZ 目录,内部读 \code{feature\_mean.npy / feature\_std.npy / feature\_cols.json},\textbf{用训练期统计量标准化推理特征}。

\subsection{STEP 7 · \code{sample\_stage2\_gflownet}}

\textbf{两种模式}:

\paragraph{(a) \code{composition\_biased}(默认)}
\begin{lstlisting}[language=bash]
python pipeline/src/19_sample_stage2_gflownet_composition_constrained.py \
  --input_dir <stage2_npz_dir> \
  --output_dir <work_dir>/stage2_gflownet_candidates_composition_decoding \
  --ckpt_path <gflownet_ckpt> \
  --split test --batch_size 128 --n_samples 100 \
  --temperature 1.0 --top_k 20 \
  --use_greedy_as_first \
  --composition_constrained \
  --target_hit_bonus 6.0 \
  --extra_element_penalty 1.0 \
  --no_overlap_penalty 6.0 \
  --stop_bias -2.0 \
  --ignore_elements H,O \
  --device <device>
\end{lstlisting}

GFlowNet 在解码每一步加 element-aware bias:
\begin{itemize}
    \item \param{target\_hit\_bonus=6}:动作引入新的目标元素 $\to$ +6 logit
    \item \param{extra\_element\_penalty=1}:动作引入非目标元素 $\to$ -1 logit
    \item \param{no\_overlap\_penalty=6}:当前 set $\cup$ \{action\} 与目标元素无任何交集 $\to$ -6 logit
    \item \param{stop\_bias=-2}:STOP 动作 logit -2(让模型生成更长的 set)
    \item \param{ignore\_elements=H,O}:H/O 不算"目标元素"
\end{itemize}

\paragraph{(b) \code{standard}} 用 \file{06\_sample\_stage2\_gflownet\_infer.py},纯模型采样,无 element bias。

\subsection{STEP 8 · \code{constrain\_stage2\_by\_composition}}

\begin{lstlisting}[language=bash]
python pipeline/src/18_constrain_stage2_sample_candidates_by_composition.py \
  --input_csv <test_samples.csv> \
  --output_csv <test_samples_composition_constrained.csv> \
  --min_coverage 0.0 --coverage_weight 20.0 \
  --extra_penalty_weight 5.0 --rank_weight 0.01 \
  --top_n_per_sample 100 \
  --dedup [--drop_zero_overlap]
\end{lstlisting}

\subsection{STEP 9 · \code{summarize\_stage2}}

聚合 step 7+8 的候选,按 set\_key 去重,统计:
\begin{itemize}
    \item \code{count}:同一 set 在所有采样中出现几次
    \item \code{frequency = count / n\_samples}
    \item \code{sample\_rank\_min/mean/max}:出现位次的统计
    \item \code{decode\_methods\_seen}:greedy / sampled / composition\_biased 等
\end{itemize}

\subsection{STEP 10 · \code{add\_composition\_fallback}}

如果一个目标的 GFlowNet 候选集合元素覆盖普遍偏低(例如目标含 P,但所有候选都没有磷源),\code{11\_add\_composition\_fallback} 会按目标元素族\textbf{直接拼}一些"标准前驱物 set":
\begin{itemize}
    \item 含 P $\to$ 加 \code{[P2O5, NH4H2PO4]}
    \item 含 S $\to$ 加 \code{[S, Na2SO4]}
    \item 含 K $\to$ 加 \code{[K2CO3, KNO3]}
    \item 含 Li $\to$ 加 \code{[Li2CO3, LiOH·H2O]}
\end{itemize}

\subsection{STEP 11 · \code{retrieve\_stage2\_candidates}}

\textbf{Case-based reasoning}:
\begin{enumerate}
    \item 在训练 NPZ 里取每个目标的 \code{(x\_struct, y\_set)}
    \item 推理目标的 \code{x\_struct}(标准化空间)对训练库做余弦相似度
    \item 返回 top-50 相似目标的真实前驱物集合(以 multi-hot label,阈值 0.5 二值化)
\end{enumerate}

\begin{repobox}[可重复要点 \#4]
Retrieval 不是"和当前目标一模一样的训练样本"——它是"特征空间最近的训练样本",返回它们的真实前驱物作为候选。这等价于"如果新目标长得像 X 化合物,就把 X 的合成路线作为参考"。
\end{repobox}

\subsection{STEP 12 · \code{predict\_stage2\_baseline}}

\textbf{多标签 ExtraTrees}:每个前驱物标签独立训练一个二分类器,推理出每个标签的概率 $P(\text{label}_i \mid x)$。

候选生成:
\begin{enumerate}
    \item 取概率 $\geq 0.02$ 的标签 top-12
    \item 按 \code{top\_k\_sets=30} 枚举 size $\leq 4$ 的所有子集组合,按 $\sum\log(P_i) + \sum\log(1-P_j)$ 打分
    \item 输出 top-30 个 set
\end{enumerate}

\textbf{降级行为}:若 model 或 script 不存在,写一个空 csv(只有表头),让下游 \code{merge\_stage2\_sources} 仍能 concat 不出错。

\subsection{STEP 13 · \code{merge\_stage2\_sources}}

四源合并(GFlowNet 已经在 fallback\_csv 里):
\begin{enumerate}
    \item \textbf{GFlowNet 多温度}(主多样性源)
    \item \textbf{Retrieval k-NN}(case-based,稳)
    \item \textbf{Composition fallback}(规则保险)
    \item \textbf{ExtraTrees baseline}(监督学习独立)
\end{enumerate}

\textbf{核心思想}:任何一源给出合理候选,后面 Stage35 都能挑出来。这是 SynPred 在 OOD 目标上不至于完全崩盘的关键工程冗余。

\subsection{STEP 14 · \code{rerank\_stage2\_by\_elements}}

公式:
\begin{equation}
\text{score}(\mathcal P, \text{target}) = 20\cdot \text{cov}(\mathcal P, \text{target}) - 5\cdot |\text{extra metals}| - 0.01\cdot \text{rank}_{\text{src}}
\end{equation}

\begin{repobox}[可重复要点 \#5]
覆盖权重 20 vs penalty 5 vs rank 0.01 比例失衡是有意的——\textbf{化学正确性 $\gg$ 排序新鲜度}。一个 cov=0.5 的候选(score = 10)永远赢一个 cov=0.0 但原 rank=1 的(score = $-0.01$)。
\end{repobox}

\subsection{STEP 15 · \code{fix\_stage2\_global\_rank}}

合并去重后,\code{rank} 列可能不连续(被去掉的 set 留 gap)。这一步把 rank 重赋为 $[1, 2, \ldots, n]$。

% --------------------------------------------------------------- §4
\section{STEP 16--19:Stage3 条件预测}

\subsection{STEP 16 · \code{build\_stage3\_features}}

同 STEP 5,但 \code{--task stage3}。\file{05\_build\_hybrid\_features\_infer.py} 内部按 \code{--task stage2|stage3} 选不同 column 子集。

\subsection{STEP 17 · \code{build\_stage3\_conditioned\_table}}

\begin{lstlisting}[language=bash]
python pipeline/src/05c_build_stage3_conditioned_feature_table_infer.py \
  --infer_hybrid_csv      <stage3_train_hybrid.csv> \
  --stage2_candidates_csv <stage2_final_csv> \
  --schema_json           <stage3_schema_json> \
  --output_csv <work_dir>/stage3_conditioned_x_*.csv \
  --max_stage2_candidates 30
\end{lstlisting}

\textbf{展开}:每个 sample $\times$ 每个 Stage2 候选 $\to$ 一行。文件大小放大 30$\times$。

\subsection{STEP 18 · \code{run\_stage3\_flow}(默认关)}

\begin{lstlisting}[language=bash]
python pipeline/src/13_run_stage3_infer_mixture_flow_conditioned.py \
  --conditioned_x_csv ... --schema_json ... \
  --flow_ckpt <ckpt> --flow_script <model.py> \
  --output_dir <out_dir>/stage3_condition_predictions_flow_*/ \
  --top_k_conditions 5 --n_flow_samples 64 \
  --device <device>
\end{lstlisting}

输出 \file{test\_candidates\_flat.csv}。

\subsection{STEP 18b · \code{run\_stage3\_lgbm}(默认开,主力)}

\begin{lstlisting}[language=bash]
python pipeline/src/13b_run_stage3_infer_lgbm_quantile.py \
  --conditioned_x_csv ... --schema_json ... \
  --temp_model_dir <runs/stage3/lgbm_quantile_ensemble_v2_fulldata/temperature> \
  --time_model_dir <runs/stage3/lgbm_quantile_ensemble_v2_fulldata/time> \
  --atm_model         <model_atmosphere_binary_final.txt> \
  --time_bucket_model <model_time_bucket.txt> \
  --output_dir <out_dir>/stage3_condition_predictions_lgbm/ \
  --top_k_conditions 5
\end{lstlisting}

\subsection{STEP 19 · \code{compare\_stage3\_models}}

如果同时跑了 flow + lgbm,算两份预测的 mean / std / min / max,按 precursor\_set group 算 top-1 温度的 absolute diff,写入 \file{stage3\_model\_comparison.md}。

% --------------------------------------------------------------- §5
\section{STEP 20--26:路线汇总与 Stage35 排序}

\subsection{STEP 20 · \code{summarize\_routes}}

\begin{lstlisting}[language=bash]
python pipeline/src/12_summarize_structure_to_routes.py \
  --stage3_flat_csv <flow_flat_csv> \
  --output_dir <out_dir>/routes_*/ \
  --top_n 100
\end{lstlisting}

\subsection{STEP 21 · \code{filter\_display\_routes}}

\begin{lstlisting}[language=bash]
python pipeline/src/14_filter_synthesis_routes_for_display.py \
  --input_csv <readable.csv> \
  --output_csv <display_filtered.csv> \
  --min_temperature_c 300 --max_temperature_c 1600 \
  --min_time_h 0.1 --max_time_h 240 \
  --top_n 30 --prefer_top_component_mean
\end{lstlisting}

\subsection{STEP 22--24 · 三种 reranker}

详见 Part II。每一步在 csv 上加 \code{stage35\_<X>\_score / stage35\_<X>\_rank} 列,\textbf{非破坏}。任一缺脚本/缺模型则 \code{record\_degradation}。

\subsection{STEP 25 · \code{best\_route\_per\_precursor}}

\textbf{逻辑}:同一 \code{precursor\_set\_key} 下,所有 (温度, 时间) 候选选最佳一组。这避免最终展示给用户时出现"前驱物 a + 800°C + 12h"与"前驱物 a + 850°C + 10h"两条几乎相同的路线。

\subsection{STEP 26 · \code{export\_final\_top\_routes}}

按优先级 v21 > learned > rule > display 选最强 ranking 写到 \file{final\_top\_routes.csv}(详见 Part II)。

% --------------------------------------------------------------- §6
\section{STEP 27 · Reliability Layer(独立一层)}

\subsection{设计意图}

\begin{quote}
\itshape 推断时间不确定性应该和排序分数解耦。
\end{quote}

Reliability layer 不改 ranking,只在 final csv 上\textbf{附加列}:置信度、警告、QC、distribution support。理由:
\begin{itemize}
    \item 排序信号是"哪条更优",置信度是"我们对这个排序有多确定"
    \item 把不确定性塞进 score 会让"高 score + 低置信度"与"中 score + 高置信度"无法区分
    \item 解耦后,UI 可以分两栏展示:\code{final\_route\_rank | confidence\_level}
\end{itemize}

\code{run\_reliability\_layer}(\file{steps\_reliability.py:1927-1984})按序调度 17 个子步骤,每个看 yaml 是否 enabled。

\subsection{核心子步骤}

\subsubsection{\code{run\_precursor\_qc}}

\begin{lstlisting}[language=bash]
python pipeline/scripts/qc_route_precursors.py \
  --input_csv <final_top_routes.csv> \
  --output_csv <final_top_routes_with_precursor_qc.csv> \
  --top_n 30
\end{lstlisting}

逐个前驱物查 RDKit/SMILES/化学合理性,标注 \code{precursor\_qc\_level} $\in$ \{pass, minor\_warning, major\_warning\} 与具体警告。

\subsubsection{\code{attach\_route\_confidence}}

核心公式:
\begin{align}
\text{score} = & \;0.30\cdot \text{element} + 0.18\cdot \text{rank} + 0.17\cdot \text{cond\_model} \nonumber \\
& + 0.17\cdot \text{cond\_rule} + 0.08\cdot \text{prob} + 0.10\cdot \text{route\_score}
\end{align}

\begin{itemize}
    \item \code{element\_score} = element\_coverage $-$ min(0.6, $0.2 \times$ missing\_count) $-$ min(0.4, $0.1 \times$ extra\_penalty)
    \item \code{rank\_score} = $1 - $ precursor\_rank / max\_rank
    \item \code{cond\_model} = stage3\_score(归一化)
    \item \code{cond\_rule} = 物理合理性
    \item \code{prob} = stage35\_v2\_prob(权重小,calibration 不好)
    \item \code{route\_score\_component} = stage35\_v21\_score / 12.0(归一化压缩)
\end{itemize}

\textbf{阈值映射}:
\begin{itemize}
    \item $\geq 0.78 \to$ \code{high\_confidence} / \code{recommended}
    \item $\geq 0.58 \to$ \code{medium\_confidence} / \code{recommended\_with\_validation}
    \item $< 0.58 \to$ \code{low\_confidence} / \code{review\_required}
\end{itemize}

\textbf{硬规则修正}:
\begin{itemize}
    \item element\_coverage $< 0.5$ OR missing\_count $\geq 2 \to$ 强制 \code{low\_confidence}
    \item 含 \code{extreme\_high\_temperature\_warning} 或 \code{too\_long\_time\_warning} $\to$ high $\to$ medium
    \item 警告级别:无警告 / 含致命警告 $\to$ major / 其余 $\to$ minor
\end{itemize}

\subsubsection{\code{attach\_stage3\_condition\_reference\_support}}

把 stage3 训练数据集里"实际观测到的 (材料族, 温度, 时间)"作为参考库,推理路线查"训练库里 $\pm 15$°C 范围有没有相似材料合成?",标 \code{condition\_reference\_support\_level} $\in$ \{strong, weak, none\}。

\subsubsection{\code{add\_condition\_distribution\_confidence}}

在所有训练样本的 (温度, 时间) 二维分布上,看推理路线落在 KDE "body"(高密度区)还是 "tail"(低密度区):
\begin{itemize}
    \item body $\to$ high score
    \item tail $\to$ low score(模型在分布尾部预测不确定性大)
\end{itemize}

\subsubsection{\code{audit\_condition\_diversity}}

\textbf{审计}:同一 sample 的 30 条候选路线,温度范围多大?如果都在 $\pm 20$°C 内,Stage3 模型可能崩。这一步只统计、不修改 ranking。

\subsubsection{\code{postprocess\_confidence\_with\_precursor\_qc}}

把 \code{precursor\_qc\_level} 与 \code{route\_confidence\_level} 联动:precursor major\_warning $\to$ confidence $-0.15$,minor $\to$ $-0.05$,严重时强制 \code{review\_required}。

\subsubsection{V3 / V43 后置 reranker}

详见 Part II:
\begin{itemize}
    \item \code{build\_joint\_route\_features} + \code{apply\_v3\_joint\_route\_rerank}
    \item \code{build\_v3\_learned\_ranker\_dataset} + \code{apply\_v3\_learned\_ranker}
    \item \code{add\_v43\_route\_template\_features} + \code{apply\_v43\_template\_ranker} + \code{apply\_v43\_safe\_strict\_gate}
\end{itemize}

\subsubsection{\code{finalize\_recommended\_routes}}

\textbf{最末端}:从所有可用 final csv 中选最强,统一字段名,生成 \file{final\_recommended\_routes.csv} + \file{.md} + \file{.json}。

\subsection{\code{\_get\_current\_csv} 模式}

reliability 层各 step 的输入选择按 fallback chain。每完成一步 enrichment,\code{\_set\_current\_csv(r, output\_csv)} 更新 \code{final\_top\_routes\_current\_csv} 为本步输出。下一步从这继续。这让"步骤可任意 enable/disable + 任意删除中间产物"都不破坏链路。

% --------------------------------------------------------------- §7
\section{STEP 28 · \code{select\_final\_recommended\_routes}}

\textbf{优先级}:
\begin{enumerate}
    \item v4.3 template-aware chem-only
    \item v3 learned reranker
    \item v3 joint reranker
    \item confidence route table
    \item original final\_top\_routes
\end{enumerate}

\textbf{关键}:\textbf{不删除、不覆盖历史输出},只创建一个稳定别名 \file{final\_recommended\_routes.csv} + \file{.md} 指向当前最强可用版本。这让用户拿同一个文件名总能读到"当前最佳"路线表。

\code{r.outputs["final\_recommended\_routes\_source"]} 记录究竟用了哪一档,写入 manifest。

% --------------------------------------------------------------- §8
\section{端到端复现指南}

\subsection{环境}

\begin{lstlisting}[language=bash]
# Python 3.10+
pip install numpy pandas scikit-learn torch yaml lightgbm joblib
pip install pymatgen ase matminer chgnet     # 结构 + GNN
pip install rdkit                             # QC (precursor)
\end{lstlisting}

\subsection{单 target 跑}

\begin{lstlisting}[language=bash]
PR=/Users/wyc/SynPred
TASK=demo_poscar_test
mkdir -p $PR/data/poscar/$TASK
cp my_target.vasp $PR/data/poscar/$TASK/

cd $PR/scripts/07_infer/structure_to_synthesis_route/pipeline
python run_pipeline.py \
  --config configs/full_route_stage3.yaml \
  --infer_name $TASK \
  --project_root $PR

ls $PR/outputs/inference/$TASK/out_dir/routes_*/
# 找:final_recommended_routes.csv / .md
\end{lstlisting}

\subsection{断点续跑}

\begin{lstlisting}[language=bash]
# 假设 STEP 19 (compare_stage3_models) 跑崩,前面都好
python run_pipeline.py \
  --config configs/full_route_stage3.yaml \
  --infer_name $TASK \
  --start_from compare_stage3_models
\end{lstlisting}

\subsection{单步调试}

\begin{lstlisting}[language=bash]
python run_pipeline.py \
  --config configs/full_route_stage3.yaml \
  --infer_name $TASK \
  --only_step stage35_v21_rerank
\end{lstlisting}

\subsection{常见坑}

\begin{center}
\small
\begin{tabularx}{\textwidth}{lXX}
\toprule
现象 & 原因 & 解 \\
\midrule
preflight 失败 & ckpt 路径错 & 检查 yaml \code{stage2.gflownet\_ckpt} 等 \\
\code{KeyError: graph\_emb\_0} & \code{--embedding\_prefix} 不对 & 看 csv 列名,默认 \code{chgnet\_<i>} 需要 prefix=chgnet \\
Stage2 候选全空 & composition\_biased 太严 & 调小 \code{target\_hit\_bonus},关 \code{drop\_zero\_overlap} \\
NPZ 标准化报错 & template\_dir 与训练期不一致 & 训练 gold\_only,推理也 gold\_only \\
\code{flow\_flat\_csv} 空 & flow ckpt 缺 & \code{steps.run\_stage3\_flow=false} \\
V21 维度不匹配 & feature\_cols.json 不配套 & 替换为同 run 的 \\
reliability 抛 & route\_out\_dir 不存在 & 检查 STEP 26 是否成功 \\
内存爆 & $100 \times 5 \times 30$ 行 OOM & 调 \code{n\_flow\_samples 32}, batch\_size 64 \\
\bottomrule
\end{tabularx}
\end{center}

% --------------------------------------------------------------- §9
\section{设计哲学}

\subsection{八条工程取舍}

\begin{enumerate}
    \item \textbf{没有大一统模型,而是异质源 + 多层 reranker}:GFlowNet/Flow 给"多样性 + 覆盖率",GBDT/ExtraTrees 给"判别力 + 稳定性",规则 / 模板给"先验"。
    \item \textbf{数据划分按 DOI/material-group}:避免同一篇论文样本被切到 train/test。
    \item \textbf{Train mode 三件套}:\code{relaxed\_only / gold\_only / curriculum}。生产路径:\code{gflownet\_joint\_rerank\_hybrid\_gold\_only\_v1}。
    \item \textbf{Stage3 模型不直接预测条件值,而是 baseline + 残差 mixture}:解耦"先验均值估计"与"分布形状学习"。
    \item \textbf{元素约束多重保险}:解码期 + 后置硬过滤 + 重排公式 + GFlowNet reward,同一先验在不同位置生效。
    \item \textbf{可靠性层独立成一层}:排序分数 $\neq$ 置信度。
    \item \textbf{Sort + Annotate, never destroy}:每个 reranker / reliability step 只新增列。
    \item \textbf{配置驱动 + 模板插值 + override}:一份 yaml 配整套。
\end{enumerate}

\subsection{与"标准"机器学习 pipeline 的差异}

\begin{center}
\small
\begin{tabular}{lll}
\toprule
维度 & 普通 ML pipeline & SynPred Stage 07 \\
\midrule
模型数 & 1 个端到端 & 6+ 个异质模型 \\
失败模式 & 一处崩全崩 & 任何一处崩 $\to$ degrade \\
排序信号 & 一个 score & rule, learned, v21, v3, v43 五种共存 \\
可解释 & 黑盒 & 每步独立 csv \\
Reproducibility & 重跑可能不一致 & manifest 记录所有产物 \\
调试 & 改代码 + 全量重跑 & --only\_step / --start\_from 单步 \\
用户接口 & 一个数 & 多档 confidence + warnings + QC \\
\bottomrule
\end{tabular}
\end{center}

\subsection{可改进点}

\begin{enumerate}[label=(\alph*)]
    \item reliability\_layer 17 个子函数耦合在一个 1984 行文件里,可拆成 sub-pipeline。
    \item Stage2/3/35 之间的中间 csv 是 flat 表格,跨 step 重新读写。可用 parquet + 内存共享。
    \item manifest 没有 hash —— 同一份输入跑两次,manifest 比对不出"产物是否真的一致"。可加 file hash。
    \item degradation 不分级 —— "CGCNN 缺"与"Flow ckpt 缺"同样写,但影响差异大。可加 \code{severity}。
    \item 没有 stage2 candidate diversity audit。
    \item 没有 dry-run 模式,可加 \code{--dry\_run}。
\end{enumerate}

% --------------------------------------------------------------- §10
\section{结语}

Stage 07 是 SynPred 的最终交付层。28 个 step 串起 4000 行代码,把"一份 POSCAR"变成"一份带置信度 + QC + 警告的合成路线推荐表"。

整个流水线的工程哲学可以一句话概括:\textbf{让多个不完美的模型协同输出一份"足够可靠"的推荐}。生成模型给广度,监督模型给深度,规则给先验,reranker chain 给精度,reliability layer 给可信度。任何单一组件都不够好,但合起来就能用。

流水线最大的优势是\textbf{优雅降级 + 可审计 + 可断点}——任何一环失效,流程不崩;任何一个决策,都能定位到具体 csv 行;任何一次失败,都能从那一步重启。这三点让 SynPred 在生产环境真正"可用"。

Stage 07 的契约只有两条:\textbf{\code{\{infer\_name\}} 隔离的 outputs 目录 + \file{pipeline\_v3\_manifest.json} 自描述}。下游(用户、UI、Dashboard、监控)只要会读这两样东西,就能完整重建任意一次推理的所有状态、所有产物、所有降级原因。这是工程"可靠性"的根本——所有信息可见、可追溯、可还原。


\end{document}
